\documentclass[11pt,a4paper]{article}
\usepackage[T1]{fontenc}
\usepackage[utf8]{inputenc}
\usepackage[french]{babel}
\usepackage{lmodern,microtype,amssymb,amsmath,mathtools,xcolor}
\usepackage[margin=2.5cm]{geometry}
\usepackage{tikz}
% What a `tikzpicture` in a transcription may rely on, beyond TikZ itself.
%
% Read twice: by grothendieck.sty, for the PDFs, and by `scripts/tikz.mjs`,
% which compiles each picture alone into the SVG the reading views show. A
% library or a colour added here reaches both; one added anywhere else reaches
% only one of them, and the PDF and the screen then disagree.

\usetikzlibrary{decorations.pathreplacing,calc,arrows.meta,patterns}

% His felt pens, as far as a scan lets us name them. Only the colours that carry
% a distinction the text relies on are defined — on the arcs of 154 p. 53 the
% light and the dark green vertices are two different subdivisions.
\definecolor{gvertclair}{HTML}{B9C230}
\definecolor{gvertfonce}{HTML}{3E5F2B}
\definecolor{gorange}{HTML}{E0782A}
\definecolor{gbrun}{HTML}{6B4A33}
\definecolor{gbleu}{HTML}{6F97A8}

\usepackage[hidelinks]{hyperref}
\newcommand{\legende}[1]{{\centering\small\color{gray!60!black}#1\par}\medskip}
\newcommand{\cote}[2]{\par\smallskip\noindent{\small\textbf{Cote #1, page #2.}}\par\nopagebreak}
\title{Schémas et diagrammes dessinés dans les notes de Montpellier\\[2pt]\large brouillon}
\author{Michel Hua}
\date{Novembre 2026}
\begin{document}
\maketitle
\begin{abstract}
Les notes manuscrites de Grothendieck conservées à Montpellier contiennent peu de dessins, mais certains représentent des schémas (l'anneau des entiers, des familles au-dessus d'une base, des courbes stables, la droite projective et ses revêtements) et d'autres sont des diagrammes commutatifs dessinés en perspective. Ce brouillon réunit les 43 dessins de ce genre que l'édition a relevés dans dix-neuf cotes, redessinés proprement, chacun avec la cote, la page et une légende qui le lit en termes actuels.
\end{abstract}

\section*{Comment lire ces figures}
Chaque dessin de la page a d'abord été redessiné fidèlement, trait pour trait, dans la transcription de sa cote : c'est ce dessin qui répond à la question de savoir ce qu'il y a sur la page. Les figures de ce brouillon en sont une seconde version, redessinée proprement d'après la première : courbes lissées, sommets alignés, croisements lisibles. Elles gardent les objets, les flèches et les étiquettes de la page ; une étiquette douteuse ou illisible le reste, et la légende le dit. Une seule forme a été corrigée, celle de la fibre nodale de la page 175 de la cote 34, pour que la section ne la coupe pas deux fois. Les légendes reprennent celles des lectures modernisées des cotes ; elles n'ont pas été vérifiées par un humain contre les fac-similés.

Ces dessins appartiennent à des manuscrits inédits, et les reproduire, même redessinés, relève des ayants droit. Ce brouillon n'est pas destiné à être déposé tel quel.


\section{Au-dessus de $\operatorname{Spec}\mathbf Z$}

Trois pages montrent l'anneau des entiers et ce qui vit au-dessus de lui comme une figure de géométrie : $\operatorname{Spec}\mathbf Z$ est une droite dont les points fermés sont les nombres premiers. À la page 173 de la cote 15, un point fermé $x$ sur cette droite sert à lire une suite de localisation. Dans la cote 63, consacrée aux courbes elliptiques, les lieux $\Delta = 0$, $j = 0$ et $j = 12^3$ sont tracés au-dessus de $\operatorname{Spec}\mathbf Z$ et se rencontrent au-dessus de $2$ et de $3$, puis, en $p = 2$, les quatre branches du niveau~$3$ se coupent en un point $a$. Ce dernier dessin est l'objet d'une section de l'article Lean (cote 63 : les quatre composantes de $j = 0$ ont quatre tangentes en $p = 2$) : Grothendieck y trace des branches distinctes, puis conclut page~99 à une tangente commune, ce que la formalisation dément.

\cote{15}{173}
\begin{center}
\begin{tikzpicture}[x=1cm,y=1cm,line cap=round,
  pt/.style={circle,fill,inner sep=1.4pt}]
\draw[thick] (0,0) -- (8,0) node[right] {$\operatorname{Spec}\mathbb{Z}$};
\node[pt,label={below:$x$}] at (3.2,0) {};
\draw[gray] (7.15,-0.2) -- (7.3,0.2);
\end{tikzpicture}
\end{center}
\legende{Figure de l'édition, redessinée d'après le dessin de la page 173 (dessin 15/09-1). $X = \operatorname{Spec}\mathbb{Z}$ avec un point fermé $x$, tel ceux de $S = X - U$ dans la suite de localisation ; le petit trait oblique près de l'extrémité droite est sur la page, sans explication.}

\cote{63}{95}
\begin{center}
\begin{tikzpicture}[x=1cm,y=1cm,line cap=round,
  pt/.style={circle,fill,inner sep=1.3pt},
  leg/.style={font=\scriptsize,align=center,inner sep=1pt}]
% la base et les deux points
\draw[thick] (0,0) -- (8.4,0) node[right] {$\operatorname{Spec}\mathbf{Z}$};
\node[pt,label={below:$p=2$}] (P2) at (2,0) {};
\node[pt,label={below:$p=3$}] (P3) at (5.5,0) {};
% le diviseur à l'infini
\draw (0,5.4) -- (8.4,5.4) node[right] {$\Delta = 0$};
% j = 0 et j = 12^3
\draw (0,3) -- (8.4,3) node[right] {$j = 0$ ($c_4 = 0$)};
\draw (0,4.1) .. controls (0.8,3.8) and (1.4,3.36) .. (2,3)
  .. controls (2.6,2.64) and (3.2,2.25) .. (3.75,2.25)
  .. controls (4.3,2.25) and (4.9,2.64) .. (5.5,3)
  .. controls (6.4,3.54) and (7.4,4.1) .. (8.4,4.3)
  node[right] {$j = 12^3$ ($c_6 = 0$)};
\node[pt] (A2) at (2,3) {};
\node[pt] (A3) at (5.5,3) {};
\draw[dashed] (A2) -- (P2);
\draw[dashed] (A3) -- (P3);
% les groupes d'automorphismes aux deux points
\node[leg] (G24) at (2.6,4.6) {groupe d'autom.\\ d'ordre $24$};
\node[leg] (G12) at (4.9,4.6) {groupe d'autom.\\ d'ordre $12$};
\draw[-{Stealth[length=1.8mm]}] (G24.south) .. controls +(0,-0.4) and +(0.25,0.4) .. (A2);
\draw[-{Stealth[length=1.8mm]}] (G12.south) .. controls +(0,-0.4) and +(-0.25,0.4) .. (A3);
% les multiplicités
\draw (3.75,2.25) -- (3.75,1.95);
\node[leg,anchor=north] at (3.75,1.9) {diviseur de\\ multiplicité $2$\\ (groupe d'autom.\\ d'ordre $4$ en\\ car.\ $\neq 2,3$)};
\draw (7.2,3) -- (7.2,2.7);
\node[leg,anchor=north] at (7.2,2.65) {diviseur de\\ multiplicité $3$\\ (groupe d'autom.\\ d'ordre $6$ en\\ car.\ $\neq 2,3$)};
\end{tikzpicture}
\end{center}
\legende{Figure de l'édition, redessinée d'après le dessin de la page 95 (dessin 63/05-1). Au-dessus de $\operatorname{Spec}\mathbf{Z}$, le diviseur à l'infini $M^{\infty} = \operatorname{div}\Delta$ et les diviseurs $M' = V(c_4)$ ($j = 0$) et $M'' = V(c_6)$ ($j = 12^3$), qui se rencontrent au-dessus de $p = 2$ et de $p = 3$ ; les légendes sont celles de la page.}

\cote{63}{99}
\begin{center}
\begin{tikzpicture}[x=1cm,y=1cm,line cap=round,
  pt/.style={circle,fill,inner sep=1.3pt}]
% (a) page 97
\begin{scope}
\foreach \a in {-43,-14,15,42} \draw (\a:1.35) -- (\a+180:1.35);
\node[pt] (a) at (0,0) {};
\node[above=2pt] at (0,0.05) {$a$};
\draw[decorate,decoration={brace,amplitude=5pt}] (-1.55,-1.0) -- (-1.55,1.0)
  node[midway,left=6pt] {$M_3$};
\node (Mp) at (2.75,0.15) {$M'_3$};
\foreach \a in {42,15,-14}
  \draw[-{Stealth[length=1.8mm]}] (Mp.west) .. controls +(-0.3,0) and ($(\a:1.15)+(0.3,0)$) .. ($(\a:1.15)+(0.08,0)$);
\draw[densely dotted] (a) -- (0,-1.8);
\draw[very thick] (-1.5,-1.8) -- (1.6,-1.8) node[pos=0,left] {$\mu_3^{*}$};
\node[pt] (b) at (0,-1.8) {};
\node[anchor=north west,inner sep=2pt] at (0.05,-1.85) {$b = \operatorname{Spec}\mathbf{F}_4$};
\draw (-1.5,-3.0) -- (1.6,-3.0) node[pos=0,left] {$\operatorname{Spec}\mathbf{Z}$};
\node[pt,label={below:$\operatorname{Spec}\mathbf{F}_2$}] at (0,-3.0) {};
\node at (0,-4.1) {(a)};
\end{scope}
% (b) page 99
\begin{scope}[xshift=7cm]
\foreach \a in {-12,-4,5,14} \draw (\a:1.45) -- (\a+180:1.45);
\draw (0,-0.6) -- (0,0.6);
\node[pt] (a') at (0,0) {};
\node[anchor=north west,inner sep=1pt] at (0.05,-0.12) {$a$};
\draw[decorate,decoration={brace,amplitude=5pt}] (-1.6,-0.5) -- (-1.6,0.5)
  node[midway,left=6pt] {$M'_3$};
\draw[decorate,decoration={brace,amplitude=6pt,mirror}] (1.75,-0.9) -- (1.75,0.9)
  node[midway,right=7pt] (M3) {$M_3$};
\draw[very thick] (-1.5,-1.8) -- (1.6,-1.8) node[pos=0,left] {$\mu_3^{*}$};
\node[pt,label={above right:$b$}] at (0,-1.8) {};
\draw (-1.5,-3.0) -- (1.6,-3.0) node[pos=0,left] {$\mathbf{Z}$};
\node[pt,label={below:$2$}] at (0,-3.0) {};
\draw[-{Stealth[length=1.8mm]}] (2.3,-0.95) -- (2.3,-1.6);
\draw[-{Stealth[length=1.8mm]}] (2.3,-2.05) -- (2.3,-2.8);
\node at (0,-4.1) {(b)};
\end{scope}
\end{tikzpicture}
\end{center}
\legende{Figure de l'édition, redessinée d'après les dessins des pages 97 (a) et 99 (b) (dessins 63/05-2 et 63/05-3 ; le croquis biffé de la page 99 est omis). Au-dessus du point $b = \operatorname{Spec}\mathbf{F}_4$ de $\mu_3^{*}$, lui-même au-dessus de $2$, les branches de $M'_3 = M_3 \cap V(c_4)$ dans $M_3$ passent toutes par le point $a$. Les branches y sont dessinées comme des droites distinctes passant par $a$ ; la page 99 conclut pourtant qu'elles ont même tangente (voir ci-dessous).}


\section{Familles au-dessus d'une base}

Une famille de variétés au-dessus d'une base $S$ se dessine comme un paquet de fibres au-dessus d'une ligne. Les croquis de cette partie suivent cette convention : un faisceau de droites au-dessus de $\operatorname{Spec} W$ dans les notes sur les cristaux (cote 7), une base balayée de courbes dans les notes sur la bonne réduction (cote 100), et, dans les notes de cohomologie étale de la cote 34, la fibre spéciale $X_s$ d'une courbe relative et ses sections $g_i(S)$ : une section la coupe en un point, deux sections ne passent pas par un même point, aucune ne passe par un point double. Ces deux interdits sont les conditions du paragraphe 2.3 de la cote 34, et la figure les rend visibles.

\cote{7}{47}
\begin{center}
\begin{tikzpicture}[x=0.75cm,y=0.75cm,line cap=round,line join=round,
  pt/.style={circle,fill,inner sep=1.3pt}]
% le cube Y (arêtes visibles seulement, comme sur la page)
\begin{scope}[shift={(-0.8,3.4)}]
\draw (0,0) rectangle (1.6,1.6);
\draw (0,1.6) -- (0.7,2.1) -- (2.3,2.1) -- (1.6,1.6);
\draw (2.3,2.1) -- (2.3,0.5) -- (1.6,0);
\node at (0.8,0.8) {$Y$};
\end{scope}
% le carré M, sa fibre verticale et le faisceau de droites
\draw (-1.4,-1.1) rectangle (1.6,1.9);
\draw (0,-1.1) -- (0,1.9);
\node[pt] (t) at (0,0.4) {};
\foreach \a in {168,205,20,2,-25} \draw (t) -- ++(\a:1.05);
\node at (2.0,0.4) {$M$};
% la base S
\draw[thick] (-1.4,-3.0) -- (1.6,-3.0) node[right] {$\operatorname{Spec} W = S$};
\node[pt] at (0,-3.0) {};
% les flèches
\draw[-{Stealth[length=2mm]}] (0.35,3.25) -- (0.35,2.05) node[midway,right] {$f$};
\draw[-{Stealth[length=2mm]}] (-1.75,-0.9) -- (-1.75,-2.8) node[midway,left] {$h$};
\end{tikzpicture}
\end{center}
\legende{Figure de l'édition, redessinée d'après le dessin de la page 47 (dessin 7/03-1). Le schéma $Y$, propre et lisse sur $M$, au-dessus de $M$, lui-même au-dessus de $S = \operatorname{Spec} W$ ; on y lit, semble-t-il, le point $t_0$ sur la fibre de $M$ au-dessus du point fermé $s_0$ de $S$, et des sections $g$ passant par $t_0$. La page ne nomme ni les points ni les droites.}

\cote{31}{55}
\begin{center}
\begin{tikzpicture}[x=1cm,y=1cm,line cap=round,
  pt/.style={circle,fill,inner sep=1.8pt}]
\draw (0,0) -- (7,0);
\draw[thick] (1.5,-0.9) .. controls (1.75,-0.4) and (1.85,-0.2) .. (2,0)
  .. controls (2.6,0.85) and (3.4,1.45) .. (4.1,1.45)
  .. controls (4.5,1.45) and (4.75,1.3) .. (4.95,1.1);
\node[pt] at (2,0) {};
\end{tikzpicture}
\end{center}
\legende{Figure de l'édition, redessinée d'après le dessin de la page 55 (dessin 31/03-2). Une droite et une courbe qui la traverse en un point marqué ; la page n'y met aucune lettre. Le croquis précède, sur la page, le passage à deux courbes $Y_1$ et $Y_2$ se coupant en un point, celui de l'exemple 3.}

\cote{34}{170}
\begin{center}
\begin{tikzpicture}[x=0.5mm,y=0.5mm,line cap=round]
\node[anchor=east] at (8,-44) {$X = {}$};
\draw[semithick] plot[smooth, tension=0.6] coordinates {(12,-44) (20,-36) (35,-32) (48,-37) (56,-46) (60,-57) (54,-67) (43,-69) (35,-62) (36,-52) (43,-44) (51,-38) (65,-34) (90,-33)};
\end{tikzpicture}
\end{center}
\legende{Figure de l'édition, redessinée d'après le dessin de la page 170 (dessin 34/09-4). La courbe nodale du brouillon : une courbe à un point double, qui fait une boucle.}

\cote{34}{170}
\begin{center}
\begin{tikzpicture}[x=1cm,y=1cm,line cap=round,>={Stealth[length=2mm]},
  every node/.style={inner sep=2pt},
  over/.style={preaction={draw=white,line width=3.2pt,arrows=-}}]
% back: X <- X_K <- X_U~ over S <-' U <- U-bar
\node (X)  at (0,3)   {$X$};
\node (XK) at (2,3)   {$X_K$};
\node (XU) at (5,3)   {$X_{\tilde U}$};
\node (S)  at (0,0.2) {$S$};
\node (U)  at (2,0.2) {$U$};
\node (Ub) at (5,0.2) {$\bar U$};
% front: bold X <= X~ over s <- s-bar
\node (Xb) at (3.2,1.8)  {$\mathbf{X}$};
\node (Xt) at (6.2,1.8)  {$\tilde X$};
\node      at (7.35,1.8) {$F_{\bar a}$};
\node (s)  at (3.2,-1.0) {$s$};
\node (sb) at (6.2,-1.0) {$\bar s$};
% back edges
\draw[->] (XK) -- (X);
\draw[->] (XU) -- (XK);
\draw[->] ([xshift=2pt]X.south) -- ([xshift=2pt]S.north);
\draw[->] ([xshift=-3pt]S.north) to[bend left=40] ([xshift=-3pt]X.south);
\draw[->] ([xshift=2pt]XK.south) -- ([xshift=2pt]U.north);
\draw[->] ([xshift=-3pt]U.north) to[bend left=40] ([xshift=-3pt]XK.south);
\draw[{Hooks[right]}->] (U) -- (S);
\draw[->] (Ub) -- (U);
\draw (XU) -- (5,1.3);
% obliques, front to back
\draw[->] (Xb) -- (XK);
\draw[->] (Xt) -- (XU);
\draw[->] (s) -- (U);
\draw[->] (sb) -- (Ub);
% front edges, over the back ones
\draw[white,line width=5pt,shorten <=6pt,shorten >=6pt] (Xt) -- (Xb);
\draw[-{Implies},double equal sign distance,shorten <=1pt,shorten >=1pt] (Xt) -- (Xb);
\draw[white,line width=4pt,shorten <=8pt,shorten >=8pt] (Xb) -- (s);
\draw (Xb) -- (s);
\draw (Xt) -- (sb);
\draw[->] (sb) -- (s);
\end{tikzpicture}
\end{center}
\legende{Figure de l'édition, redessinée d'après le dessin de la page 170 (dessin 34/09-5), sur le brouillon que la rédaction n'intègre pas. Au fond, $X \leftarrow X_K \leftarrow X_{\tilde U}$ au-dessus de $S \hookleftarrow U \leftarrow \bar U$, avec une projection et une section entre $X$ et $S$ et entre $X_K$ et $U$ ; devant, $\mathbf{X} \Leftarrow \tilde X$ au-dessus de $s \leftarrow \bar s$, relié au fond par des flèches obliques ; $F_{\bar a}$ est écrit à droite. Le trait vertical sous $X_{\tilde U}$ s'arrête sur la page peu après avoir croisé $\mathbf{X} \Leftarrow \tilde X$ ; on le laisse inachevé.}

\cote{34}{175}
\begin{center}
\begin{tikzpicture}[x=1cm,y=1cm,line cap=round,
  pt/.style={circle,fill,inner sep=1.1pt}]
% (a) the fibre X_s cut by the section g_i(S)
\begin{scope}
\draw[semithick] (0.15,1.5) .. controls (0.05,0.6) and (-0.05,-0.6) .. (-0.2,-1.5);
\node[above] at (0.15,1.5) {$X_s$};
\draw[semithick] (-1.3,-0.15) .. controls (-0.5,0.05) and (0.6,0.05) .. (1.4,-0.05);
\node[right] at (1.4,-0.05) {$g_i(S)$};
\node[circle,draw,inner sep=1.6pt] at (0.02,0.0) {};
\node at (0,-2.0) {(a)};
\end{scope}
% (b) two sections through one point of X_s
\begin{scope}[xshift=4.3cm]
\draw[semithick] (0,1.5) .. controls (0,0.6) and (0.05,-0.6) .. (0.2,-1.5);
\node[above] at (0,1.5) {$X_s$};
\draw[semithick] (-1.3,-0.55) -- (1.3,0.55) node[right] {$g_i(S)$};
\draw[semithick] (-1.3,0.25) -- (1.3,-0.25) node[right] {$g_j(S)$};
\node[pt] at (0,0) {};
\node at (0,-2.0) {(b)};
\end{scope}
% (c) a section through the double point of X_s
\begin{scope}[xshift=8.6cm]
\draw[semithick] (-0.55,-1.5) .. controls (-0.3,-0.75) and (-0.15,-0.25) .. (0,0)
  .. controls (0.35,0.6) and (0.75,1.35) .. (0,1.45)
  .. controls (-0.75,1.35) and (-0.35,0.6) .. (0,0)
  .. controls (0.15,-0.25) and (0.3,-0.75) .. (0.55,-1.5);
\node[above] at (0,1.45) {$X_s$};
\draw[thick] (-1.3,0.1) .. controls (-0.5,0.05) and (0.5,-0.05) .. (1.3,-0.15);
\node[right] at (1.3,-0.15) {$g_i(s)$};
\node[pt] at (0,0) {};
\node at (0,-2.0) {(c)};
\end{scope}
\end{tikzpicture}
\end{center}
\legende{Figure de l'édition, redessinée d'après les trois dessins de la page 175 (dessins 34/09-6, 09-7 et 09-8), réunis en trois panneaux. (a) La section $g_i(S)$ coupe la fibre spéciale $X_s$ en un point. (b) Deux sections $g_i(S)$, $g_j(S)$ passant par un même point de $X_s$ : cas exclu, trois branches ne pouvant se croiser normalement en dimension $2$. (c) Une section passant par un point double de $X_s$ : cas exclu, les $g_i(S)$ étant contenus dans l'ouvert de lissité. Le dernier dessin écrit $g_i(s)$.}

\cote{100}{20}
\begin{center}
\begin{tikzpicture}[x=0.95mm,y=0.95mm,line cap=round,line join=round,
  >={Stealth[length=2mm]}, sweep/.style={gray,thin}]
% the region of the base
\draw[semithick] plot[smooth, tension=0.6] coordinates {(30.6,-24) (30,-40) (29.6,-57.5)};
\draw[semithick] plot[smooth, tension=0.6] coordinates {(32.7,-18) (50,-15.5) (75,-12.7) (98.8,-11.5)};
\draw[semithick] plot[smooth, tension=0.6] coordinates {(98.8,-11.5) (98.7,-28) (98.5,-40) (96,-41.6)};
\draw[semithick] plot[smooth, tension=0.6] coordinates {(96,-41.6) (88,-42.1) (80,-42.3)};
\draw[semithick] plot[smooth, tension=0.6] coordinates {(29.6,-57.5) (38.4,-57.3) (47.7,-56.2) (56.2,-55.3) (65.5,-53.4) (70.8,-50.5)};
% the six parallel curves sweeping it
\draw[sweep] plot[smooth, tension=0.6] coordinates {(29.6,-57.5) (33.1,-49.2) (38.3,-39.9) (43.4,-32.2) (48.6,-25) (53.7,-20.3) (58.6,-16.2)};
\draw[sweep] plot[smooth, tension=0.6] coordinates {(38.4,-57.3) (43.4,-47.1) (47,-38.9) (50.1,-34.7) (54.7,-30.6) (61.4,-25) (67.1,-19.8) (73,-15.2)};
\draw[sweep] plot[smooth, tension=0.6] coordinates {(47.8,-55.9) (52.2,-46.6) (56.8,-38.9) (61.4,-32.7) (67.1,-27) (74.3,-21.4) (88.2,-13.9)};
\draw[sweep] plot[smooth, tension=0.6] coordinates {(56.2,-55.4) (59.9,-48.1) (64,-42) (69.2,-36.8) (75.3,-31.1) (84.6,-25) (98.5,-14.9)};
\draw[sweep] plot[smooth, tension=0.6] coordinates {(65.8,-53.5) (70.2,-47.1) (73.3,-42.5) (77.4,-37.8) (84.6,-32.2) (92.3,-27.5) (99.2,-22.9)};
\draw[sweep] plot[smooth, tension=0.6] coordinates {(70.8,-50.5) (74.3,-47.1) (78.4,-42) (84.6,-36.8) (92.3,-31.7) (98.2,-28.1)};
% the four transversal strokes
\draw[sweep] plot[smooth, tension=0.6] coordinates {(48.5,-25.8) (50.4,-34.6) (53.8,-43) (56.5,-44.6)};
\draw[sweep] plot[smooth, tension=0.6] coordinates {(68.6,-19.1) (71.7,-21.9) (74.3,-28.6) (77.4,-34.7) (79.4,-37.8)};
\draw[sweep] plot[smooth, tension=0.6] coordinates {(67.4,-39.6) (69.2,-44) (70.7,-49.2) (70.8,-50.5)};
\draw[sweep] plot[smooth, tension=0.6] coordinates {(90.8,-22.3) (92.3,-28.5) (96.2,-38.1)};
% the marks
\node[fill=white,inner sep=1pt] at (46.6,-40.6) {$S$};
\node[fill=white,inner sep=1pt] at (56.2,-41.6) {$S_0$};
\node[fill=white,inner sep=1pt] at (73.6,-41) {$S_1$};
\node[fill=white,inner sep=1pt] (Sp) at (88.6,-39.6) {$S'$};
\draw[->] (Sp) -- (77.6,-41.2);
\node[fill=white,inner sep=1pt] (A) at (80.4,-28.4) {$A_{S'}$};
\node[fill=white,inner sep=1pt] (B) at (96.8,-28.4) {$B_{S'}$};
\draw[->] (A) -- (B);
% the ringed point below the lower left corner, and the arrow ending at it
\draw (33.5,-60.6) circle[radius=2.6];
\fill (33.5,-60.6) circle[radius=0.75];
\draw[->] (23.5,-60.6) -- (30.6,-60.6);
\end{tikzpicture}
\end{center}
\legende{Figure de l'édition, redessinée d'après le dessin de la page 20 (dessin 100/01-4). La base $S$ figurée comme une région balayée de courbes parallèles (en gris, avec quatre traits transverses), sur laquelle sont marqués $S_0$ et $S_1 \leftarrow S'$, et, au-dessus de $S'$, la flèche $A_{S'} \to B_{S'}$ ; sous le coin inférieur gauche, un point cerclé où aboutit une flèche. La page ne dit pas ce que désignent $S_0$, $S_1$ et le point cerclé.}


\section{Courbes stables et graphes}

Dans les notes sur les multiplicités modulaires et la structure à l'infini des $M_{g,\nu}$ (cotes 146 et 147), Grothendieck dessine les courbes stables du bord : des droites projectives qui se coupent en des points doubles et portent des points marqués. Il place souvent en regard le graphe dual, un sommet par composante, une arête par point double, les sommets marqués de leur genre. Les pages 140 à 142 de la cote 147 passent ensuite aux voisinages tubulaires d'une courbe à point double et de ses branches.

\cote{146}{67}
\begin{center}
\begin{tikzpicture}[x=1.1mm, y=1.1mm, line cap=round,
  comp/.style={line width=0.7pt},
  pt/.style={circle, fill, inner sep=1.1pt},
  every node/.style={font=\small}]
% une droite portant deux grappes de composantes, marquée 5
\draw[comp] (14.5,-9) -- (25,-8.5);
\draw[comp] (14.5,-9) -- (13,-3) (14.5,-9) -- (16,-2.5);
\draw[comp] (25,-8.5) -- (24,-3.5) (25,-8.5) -- (27,-3) (25,-8.5) -- (28.5,-3.5);
\node at (10.5,-3.5) {$5$};
% deux droites secantes : un point 5*, i sur l'une, trois points sur l'autre
\coordinate (A1) at (18,-12); \coordinate (A2) at (31,-27);
\coordinate (B1) at (20.5,-36.5); \coordinate (B2) at (28.5,-13.5);
\draw[comp] (A1) -- (A2);
\draw[comp] (B1) -- (B2);
\node[pt] (i) at ($(A1)!0.27!(A2)$) {};
\node[left=3pt] at (i) {$5^*$};
\node[below left=0pt] at ($(i)+(-0.3,-0.6)$) {$i$};
\foreach \t in {0.30, 0.17, 0.065} \node[pt] at ($(B1)!\t!(B2)$) {};
\node at (10.5,-24) {$1 \leqslant i \leqslant 4$};
% une droite portant trois composantes, deux d'entre elles fourchues
\draw[comp] (49,-23) -- (68,-20);
\draw[comp] (49,-23) -- (48,-15.5) (49,-23) -- (51,-14);
\draw[comp] (59.5,-21.5) -- (59,-14);
\draw[comp] (68,-20) -- (67,-12.5) (68,-20) -- (72.5,-13);
% trois droites en chaine, portant des points ; la marque P
\coordinate (C1) at (45,-27.5); \coordinate (C2) at (53,-42);
\coordinate (D1) at (46,-38.5); \coordinate (D2) at (78.5,-34.5);
\coordinate (E1) at (70.5,-43.5); \coordinate (E2) at (75,-24.5);
\draw[comp] (C1) -- (C2);
\draw[comp] (D1) -- (D2);
\draw[comp] (E1) -- (E2);
\foreach \t in {0.17, 0.40} \node[pt] at ($(C1)!\t!(C2)$) {};
\node[pt] at ($(D1)!0.51!(D2)$) {};
\foreach \t in {0.76, 0.95} \node[pt] at ($(E1)!\t!(E2)$) {};
\node at (90,-23.5) {$\mathbb{P}$};
\end{tikzpicture}
\end{center}
\legende{Figure de l'édition, redessinée d'après le dessin de la page 67
(dessin 146/04-10). Les configurations que la page dessine en cherchant ce qui
manque à $\widehat{M}'_{0,5}$ : chaque composante est une droite projective,
chaque point un point marqué ; les marques $5$, $5^*$, $i$, $1 \leqslant i
\leqslant 4$ et $\mathbb{P}$ sont celles de la page, qui ne les commente pas.}

\cote{146}{73}
\begin{center}
\begin{tikzpicture}[x=0.8mm, y=0.8mm, line cap=round,
  comp/.style={line width=0.7pt}]
\draw[comp] (29,-60.5) -- (48,-60.5);
\draw[comp] (51.5,-61) -- (52.5,-26);
\draw[comp] (51.5,-61) -- (62.5,-29.5);
\draw[comp] (51.5,-61) -- (73,-36);
\end{tikzpicture}
\end{center}
\legende{Figure de l'édition, redessinée d'après le dessin de la page 73
(dessin 146/04-12). La seconde des deux petites configurations de courbes de
la page, la première étant biffée : une composante horizontale et un éventail
de trois composantes issues d'un même point ; la page ne permet pas de décider
si le trait horizontal appartient à la première configuration ou à la seconde,
et il est laissé à sa place.}

\cote{147}{52}
\begin{center}
\begin{tikzpicture}[scale=0.85, line cap=round,
  comp/.style={line width=0.7pt},
  pt/.style={circle, fill, inner sep=1.1pt},
  every node/.style={font=\small}]
% une droite portant a, b, c
\coordinate (L1) at (0,0.2); \coordinate (L2) at (3,1.4);
\draw[comp] (L1) -- (L2);
\foreach \t/\l in {0.22/a, 0.5/b, 0.78/c}
  \node[pt, label={[inner sep=1pt]above left:$\l$}] at ($(L1)!\t!(L2)$) {};
% deux droites secantes portant a, b et c, d
\begin{scope}[xshift=4.2cm]
\coordinate (P1) at (0,0); \coordinate (P2) at (3,1.6);
\coordinate (Q1) at (1.0,1.76); \coordinate (Q2) at (3.2,-0.44);
\draw[comp] (P1) -- (P2);
\draw[comp] (Q1) -- (Q2);
\foreach \t/\l in {0.15/a, 0.38/b}
  \node[pt, label={[inner sep=1pt]above left:$\l$}] at ($(P1)!\t!(P2)$) {};
\foreach \t/\l in {0.62/c, 0.84/d}
  \node[pt, label={[inner sep=1pt]below left:$\l$}] at ($(Q1)!\t!(Q2)$) {};
\end{scope}
% une chaine de trois droites : a, b ; c ; d, e
\begin{scope}[xshift=8.4cm]
\coordinate (M1) at (0,1); \coordinate (M2) at (4,1);
\coordinate (G1) at (0.1,-0.2); \coordinate (G2) at (1.25,1.6);
\coordinate (H1) at (2.75,1.6); \coordinate (H2) at (3.9,-0.2);
\draw[comp] (M1) -- (M2);
\draw[comp] (G1) -- (G2);
\draw[comp] (H1) -- (H2);
\node[pt, label={[inner sep=1pt]above:$c$}] at (2,1) {};
\foreach \t/\l in {0.12/a, 0.36/b}
  \node[pt, label={[inner sep=1pt]left:$\l$}] at ($(G1)!\t!(G2)$) {};
\foreach \t/\l in {0.64/d, 0.88/e}
  \node[pt, label={[inner sep=1pt]right:$\l$}] at ($(H1)!\t!(H2)$) {};
\end{scope}
\end{tikzpicture}
\end{center}
\legende{Figure de l'édition, redessinée d'après le dessin de la page 52
(dessin 147/03-1). De gauche à droite, les courbes stables des trois arbres de
(64), une droite projective par sommet et un point double par arête : la droite
à trois points ($\nu = 3$), deux droites sécantes ($\nu = 4$), trois droites en
chaîne ($\nu = 5$). La page les place l'une sous l'autre, en regard des arbres.}

\cote{147}{53}
\begin{center}
\begin{tikzpicture}[line cap=round,
  comp/.style={line width=0.7pt},
  pt/.style={circle, fill, inner sep=1.1pt},
  every node/.style={font=\small}]
% les trois graphes de (65)
\node[pt] (u) at (0.4,0) {};
\draw[comp] (u) .. controls (1.9,1.0) and (1.9,-1.0) .. (u);
\node[left] at (u) {$1$};
\node[pt] (v1) at (0,-2) {}; \node[pt] (v2) at (1.6,-2) {};
\draw[comp] (v1) to[bend left=40] (v2);
\draw[comp] (v1) to[bend right=40] (v2);
\node[left] at (v1) {$1$}; \node[right] at (v2) {$1$};
\node[pt] (w2) at (0,-4) {}; \node[pt] (w1) at (0.9,-4) {};
\draw[comp] (w2) -- (w1);
\draw[comp] (w1) .. controls (2.4,-3.0) and (2.4,-5.0) .. (w1);
\node[above] at (w2) {$2$};
% les courbes correspondantes
\begin{scope}[xshift=5cm]
% une courbe nodale portant a
\draw[comp] (-0.6,0.45) .. controls (0.3,-0.2) and (1.3,-0.6) .. (1.3,0)
  .. controls (1.3,0.6) and (0.3,0.2) .. (-0.6,-0.45);
\node[pt, label={[inner sep=1pt]right:$a$}] at (1.3,0) {};
% deux droites se coupant en deux points, portant a et b
\draw[comp] (-0.8,-2.35) .. controls (0.3,-1.2) and (1.4,-1.2) .. (2.4,-2.35);
\draw[comp] (-0.8,-1.65) .. controls (0.3,-2.8) and (1.4,-2.8) .. (2.4,-1.65);
\node[pt, label={[inner sep=1pt]above:$a$}] at (0.8,-1.49) {};
\node[pt, label={[inner sep=1pt]below:$b$}] at (0.8,-2.51) {};
% une courbe nodale a laquelle est attachee une droite portant a et b
\draw[comp] (-0.6,-3.55) .. controls (0.3,-4.2) and (1.3,-4.6) .. (1.3,-4)
  .. controls (1.3,-3.4) and (0.3,-3.8) .. (-0.75,-4.56);
\coordinate (L1) at (-0.85,-3.79); \coordinate (L2) at (0.35,-5.47);
\draw[comp] (L1) -- (L2);
\foreach \t/\l in {0.6/a, 0.88/b}
  \node[pt, label={[inner sep=3pt]left:$\l$}] at ($(L1)!\t!(L2)$) {};
\end{scope}
\end{tikzpicture}
\end{center}
\legende{Figure de l'édition, redessinée d'après le dessin de la page 53
(dessin 147/03-5). À gauche, les trois graphes à un cycle de (65), avec le
poids de leurs sommets : la boucle sur un sommet de poids $1$, deux sommets de
poids $1$ reliés par deux arêtes, la boucle munie d'une queue dont le bout porte
deux points ; à droite, en regard, les courbes stables correspondantes : la
cubique nodale portant $a$, deux droites qui se coupent en deux points, une
courbe nodale à laquelle est attachée une droite portant $a$ et $b$. Sur la
page, le $a$ de la première courbe est peu visible.}

\cote{147}{57}
\begin{center}
\begin{tikzpicture}[line cap=round,
  comp/.style={line width=0.7pt},
  pt/.style={circle, fill, inner sep=1.1pt},
  every node/.style={font=\small}]
% ligne 1 : une arete entre un sommet 3 (alpha_0) et un sommet 2 ; deux droites secantes
\node[pt] (a1) at (0,0) {}; \node[pt] (a2) at (1.4,0) {};
\draw[comp] (a1) -- (a2);
\node[left] at (a1) {$3$}; \node[below=1pt] at (a1) {$\alpha_0$};
\node[right] at (a2) {$2$};
\begin{scope}[xshift=5.6cm]
\coordinate (X) at (0,0);
\draw[comp] ($(X)+(-1.6,0.27)$) -- ($(X)+(1.2,-0.2)$);
\draw[comp] ($(X)+(-1.3,-0.55)$) -- ($(X)+(2.0,0.84)$);
\foreach \s/\l in {-1.3/a, -0.8/b}
  \node[pt, label={[inner sep=1pt]above:$\l$}] at ($(X)+(\s,-0.17*\s)$) {};
\foreach \s/\l in {0.6/c, 1.1/d, 1.6/e}
  \node[pt, label={[inner sep=1pt]above left:$\l$}] at ($(X)+(\s,0.42*\s)$) {};
\end{scope}
% ligne 2 : une boucle sur un sommet 2 (alpha_0) ; une courbe nodale portant a et b
\begin{scope}[yshift=-1.9cm]
\node[pt] (b1) at (0,0) {};
\draw[comp] (b1) .. controls (1.6,0.9) and (1.6,-0.9) .. (b1);
\node[above=1pt] at (b1) {$2$}; \node[below=1pt] at (b1) {$(\alpha_0)$};
\begin{scope}[xshift=4.8cm]
\draw[comp] (-0.6,0.45) .. controls (0.3,-0.2) and (1.3,-0.6) .. (1.3,0)
  .. controls (1.3,0.6) and (0.3,0.2) .. (-0.6,-0.45);
\node[pt, label={[inner sep=1pt]above right:$a$}] at (1.189,0.246) {};
\node[pt, label={[inner sep=1pt]below right:$b$}] at (1.189,-0.246) {};
\end{scope}
\end{scope}
% ligne 3 : une arete entre alpha_0 (g = 1) et un sommet 2 (g = 0) ; deux branches
\begin{scope}[yshift=-3.8cm]
\node[pt] (c1) at (0,0) {}; \node[pt] (c2) at (1.4,0) {};
\draw[comp] (c1) -- (c2);
\node[above=1pt] at (c2) {$2$};
\node[below=1pt] (al) at (c1) {$(\alpha_0)$};
\node[below=-1pt] at (al.south) {$g_{\alpha_0} = 1$};
\node[below=1pt] at ($(c2)+(0.35,0)$) {$g_\alpha = 0$};
\begin{scope}[xshift=5.4cm]
\coordinate (Y) at (0,0);
\draw[comp] (-1.5,-0.9) .. controls (-0.9,-0.6) and (-0.4,-0.35) .. (Y) -- (0.3,0.2);
\coordinate (Z1) at (-0.3,0.22); \coordinate (Z2) at (1.9,-1.38);
\draw[comp] (Z1) -- (Z2);
\node[circle, fill, inner sep=1.6pt] at (Y) {};
\foreach \t/\l in {0.42/a, 0.74/b}
  \node[pt, label={[inner sep=1pt]above right:$\l$}] at ($(Z1)!\t!(Z2)$) {};
\node[below=2pt] at (-1.1,-0.75) {$g_{\alpha_0} = 1$};
\end{scope}
\end{scope}
\end{tikzpicture}
\end{center}
\legende{Figure de l'édition, redessinée d'après le dessin de la page 57
(dessin 147/03-9). Les trois MD-graphes de dimension ambiante $2$ à une seule
arête, chacun en regard de la courbe stable qu'il représente : deux droites
sécantes, $\alpha_0$ portant trois points (type $(0, 5)$) ; une courbe nodale
portant deux points (type $(1, 2)$) ; une composante de genre $1$ coupant, au
point épais, une droite qui porte deux points (type $(1, 2)$). L'indice
surchargé du second sommet de la première ligne n'est pas lu et n'est pas
reporté.}

\cote{147}{140}
\begin{center}
\begin{tikzpicture}[x=0.45mm, y=0.45mm, line cap=round,
  every node/.style={font=\small}]
\def\boucle{(14.7,-33) (33.3,-38.3) (53.3,-47.3) (70,-73.3) (71.7,-90)
  (60,-110) (48.7,-114.3) (30,-106.7) (23.3,-96.7) (30,-76.7) (53.3,-47.3)
  (81.7,-29.3)}
% (a) la courbe qui se recoupe en un point double
\draw[line width=0.7pt] plot[smooth, tension=0.55] coordinates {\boucle};
\node[circle, fill, inner sep=1.4pt] at (53.3,-47.3) {};
\node at (48,-130) {(a)};
% (b) la meme boucle epaissie en une bande, qui se croise au point double
\begin{scope}[xshift=110]
\draw[line width=5.6mm, line cap=butt] plot[smooth, tension=0.55] coordinates {\boucle};
\draw[line width=5.1mm, line cap=butt, black!12] plot[smooth, tension=0.55] coordinates {\boucle};
\node[circle, fill, inner sep=1.4pt] at (53.3,-47.3) {};
\node at (48,-130) {(b)};
\end{scope}
\end{tikzpicture}
\end{center}
\legende{Figure de l'édition, redessinée d'après les deux premiers dessins de
la page 140 (dessins 147/07-1 et 147/07-2), en deux panneaux. (a) L'image dans
$Y$ d'une courbe immergée $D'$ qui se recoupe en un point double. (b) La même
boucle épaissie en une bande, dont les deux branches se croisent au point
double. Le troisième dessin de la page, une bande qui ne se recoupe pas, n'est
pas repris ici.}

\cote{147}{141}
\begin{center}
\begin{tikzpicture}[line cap=round,
  pt/.style={circle, fill, inner sep=1.1pt},
  every node/.style={font=\small}]
% le voisinage tubulaire V_{D'}, bande autour de la courbe D'
\filldraw[fill=black!10, line width=0.5pt]
  (0,0) .. controls (2.8,-1) and (5.7,-1) .. (8.5,0) -- (8.5,-1)
  .. controls (5.7,-2) and (2.8,-2) .. (0,-1) -- cycle;
% les deux disques V_{D''} autour des points de D''
\coordinate (p1) at (1.01,-0.82);
\coordinate (p2) at (5.62,-1.17);
\filldraw[fill=black!28, line width=0.5pt] (p1) circle (0.38) (p2) circle (0.38);
\draw[line width=0.8pt] (0,-0.5) .. controls (2.8,-1.5) and (5.7,-1.5) .. (8.5,-0.5)
  node[right] {$D'$};
\node[pt] at (p1) {}; \node[pt] at (p2) {};
\node at (-0.55,0.3) {$V_{D'}$};
\node (dd) at (4.3,0.45) {$D''$};
\draw[dashed, -{Stealth[length=1.8mm]}, shorten >=2.5pt] (dd) -- (p1);
\draw[dashed, -{Stealth[length=1.8mm]}, shorten >=2.5pt] (dd) -- (p2);
\node (vd) at (3.3,-2.75) {$V_{D''}$};
\draw[dashed, -{Stealth[length=1.8mm]}] (vd) -- ($(p1)!0.4cm!(vd)$);
\draw[dashed, -{Stealth[length=1.8mm]}] (vd) -- ($(p2)!0.4cm!(vd)$);
\end{tikzpicture}
\end{center}
\legende{Figure de l'édition, redessinée d'après le croquis de la page 141
(dessin 147/08-1). Le voisinage tubulaire $V_{D'}$, bande autour de la courbe
$D'$ ; deux points de $D''$ sur $D'$, chacun dans un petit disque, son
voisinage $V_{D''}$ : la situation $D'' \to D' \to Y$ du diagramme de
fonctorialité.}

\cote{147}{142}
\begin{center}
\begin{tikzpicture}[x=0.95mm, y=0.95mm, line cap=round,
  every node/.style={font=\small}]
\def\haut{(2.4,-9.8) (26.8,-18.3) (63.4,-14.6) (100,-19.5)}
\def\ame{(8.5,-18.3) (23.2,-24.2) (45.1,-25.8) (78,-27.6) (100,-25.6)}
% la bande, coupee a droite par le bord du dessin
\fill[black!10] plot[smooth, tension=0.55] coordinates {\haut}
  -- plot[smooth, tension=0.55] coordinates {(100,-34.2) (75.6,-37.8) (51.2,-39) (20.7,-34.2) (4.9,-22)}
  -- cycle;
\draw[line width=0.5pt] plot[smooth, tension=0.55] coordinates {\haut};
\draw[line width=0.5pt] plot[smooth, tension=0.55] coordinates {(4.9,-22) (20.7,-34.2) (51.2,-39) (75.6,-37.8) (100,-34.2)};
\draw[line width=0.5pt] (2.4,-9.8) -- (4.9,-22);
% autour des deux points de croisement, les quatre quadrants hachures
\foreach \p/\u/\l in {{(23.2,-24.2)}/{(25.5,-16.0)}/{(20.9,-32.4)}, {(78,-27.6)}/{(77.4,-19.6)}/{(78.6,-35.6)}} {
  \fill[pattern=north east lines] \p circle (5.5);
  \begin{scope}
    \clip \p circle (5.5);
    \draw[black!10, line width=3pt] plot[smooth, tension=0.55] coordinates {\ame};
    \draw[black!10, line width=3pt] \u -- \l;
  \end{scope}
  \draw[line width=0.5pt] \p circle (5.5);
}
% la courbe centrale et les deux branches transverses
\draw[line width=1.1pt] plot[smooth, tension=0.55] coordinates {\ame};
\draw[line width=0.8pt] (25.5,-16.0) -- (20.9,-32.4);
\draw[line width=0.8pt] (77.4,-19.6) -- (78.6,-35.6);
\node[circle, fill, inner sep=1.2pt] at (23.2,-24.2) {};
\node[circle, fill, inner sep=1.2pt] at (78,-27.6) {};
% la marque Theta_{D',Y}
\node (th) at (55,-7.5) {$\Theta_{D',Y}$};
\draw[-{Stealth[length=1.8mm]}, shorten >=1pt] (th) -- (25.3,-16.8);
\draw[-{Stealth[length=1.8mm]}, shorten >=1.5pt] (th) -- (47,-25.9);
\end{tikzpicture}
\end{center}
\legende{Figure de l'édition, redessinée d'après le dessin de la page 142
(dessin 147/08-2). Le voisinage tubulaire $V_{D',Y}$ (la bande, coupée à
droite par le bord du dessin), traversé par la courbe $D'$ ; en deux points,
une branche transverse la coupe, et autour de chacun quatre quadrants hachurés.
La marque $\Theta_{D',Y}$, image inverse de $\Theta$ dans $V_{D',Y}$,
désigne de ses flèches une branche transverse et la courbe $D'$.}

\cote{147}{180}
\begin{center}
\begin{tikzpicture}[line cap=round]
\draw[line width=0.7pt] (-1.1,0.35) .. controls (-0.2,0.2) and (0.5,-0.2) .. (0.5,-0.7)
  .. controls (0.5,-1.3) and (-0.5,-1.3) .. (-0.5,-0.7)
  .. controls (-0.5,-0.2) and (0.2,0.2) .. (1.1,0.35);
\end{tikzpicture}
\end{center}
\legende{Figure de l'édition, redessinée d'après le dessin de la page 180
(dessin 147/09-1). Une courbe qui se recoupe en un point double : la courbe
nodale, dont la stratification par son nœud est exclue par la condition d).}


\section{La droite projective et ses revêtements}

La cote 138 est faite de dessins d'enfant. La droite projective privée de $0$, $1$, $\infty$ y est vue comme une sphère triangulée par trois grands cercles, et le revêtement $x \mapsto x^n$ par ses fibres au-dessus des trois points.

\cote{138}{43}
\begin{center}
\begin{tikzpicture}[x=1cm,y=1cm,line cap=round,line join=round,
  pt/.style={circle,fill,inner sep=1.3pt},
  opt/.style={circle,draw,fill=white,inner sep=1.3pt},
  hid/.style={gray,densely dashed}]
% (a) la sphère et ses trois grands cercles
\begin{scope}[x=1.5cm,y=1.5cm]
% équateur : 0, 1, 0, 1
\draw[hid] plot[domain=0:180,samples=60,smooth] ({cos(\x)},{0.35*sin(\x)});
\draw[thick] plot[domain=180:360,samples=60,smooth] ({cos(\x)},{0.35*sin(\x)});
% grand cercle par les deux 0 et les deux infinis
\draw[thick] plot[domain=0:180,samples=60,smooth] ({cos(\x)},{0.937*sin(\x)});
\draw[hid] plot[domain=180:360,samples=60,smooth] ({cos(\x)},{0.937*sin(\x)});
% grand cercle par les deux 1 et les deux infinis
\draw[thick] plot[domain=65.5:245.5,samples=60,smooth]
  ({0.574*cos(\x)},{0.35*0.819*cos(\x)+0.937*sin(\x)});
\draw[hid] plot[domain=-114.5:65.5,samples=60,smooth]
  ({0.574*cos(\x)},{0.35*0.819*cos(\x)+0.937*sin(\x)});
% l'axe des infinis
\draw[hid] (0,0.937) -- (0,-0.937);
\draw (0,1.02) -- (0,1.3);
\draw (0,-1.02) -- (0,-1.3);
\node[pt,label={[inner sep=1pt,xshift=-2pt]left:$\infty$}] at (0,0.937) {};
\node[pt,label={[inner sep=1pt]right:$\infty$}] at (0,-0.937) {};
\node[opt,label={left:$0$}] at (-1,0) {};
\node[opt,label={right:$0$}] at (1,0) {};
\node[pt,label={[inner sep=1pt]above right:$1$}] at (0.574,0.287) {};
\node[pt,label={[inner sep=1pt]below left:$1$}] at (-0.574,-0.287) {};
\node[pt] at (0.174,-0.345) {};
\node at (0,-1.7) {(a)};
\end{scope}
% (b) la seconde figure, partielle
\begin{scope}[shift={(2.5,7.1)},x=0.5mm,y=0.5mm]
\draw[thick] plot[smooth cycle, tension=0.6] coordinates {(8.5,-166) (6.5,-156.8)
  (7.2,-143.8) (12.4,-130.7) (20.3,-121.5) (29.4,-116.3) (38.6,-114.4)
  (43.8,-116.3) (49,-122.8) (55.6,-130.7) (59.5,-137.2) (60.8,-146.4)
  (59.5,-154.2) (57.6,-159.4) (56.9,-165.3) (53,-168.6) (45.1,-169.9)
  (38.6,-168.2) (28.1,-169.3) (19,-168.6) (13.7,-169.3)};
\draw[thick] plot[smooth cycle, tension=0.6] coordinates {(44.1,-132.9)
  (37.3,-132.2) (28.1,-133.9) (20.3,-138.5) (16.4,-143.8) (15.7,-149)
  (17.7,-154.2) (22.9,-158.1) (29.7,-160.1) (37.3,-160.7) (47.7,-158.8)
  (55.6,-154.8) (58.8,-149.6) (59.5,-143.8) (56.9,-137.9) (50.3,-133.9)};
\draw[gray] plot[smooth, tension=0.6] coordinates {(49,-135) (53.5,-138)
  (56.3,-141.5) (57.3,-144.5)};
\draw (12.4,-146) -- (19.4,-146);
\draw (56.4,-149.6) -- (62.4,-149.6);
\node[left,inner sep=1pt] at (12.6,-146) {$1$};
\node[right,inner sep=1pt] at (62.2,-149.6) {$1$};
\node[pt,label={[inner sep=1pt]above:$\infty$}] at (38.6,-114.4) {};
\node[pt,label={[inner sep=1pt]below:$\infty$}] at (38.6,-168.2) {};
\node[pt,label={[inner sep=1pt]below left:$0$}] at (44.1,-132.9) {};
\node[pt,label={[inner sep=1pt]above left:$0$}] at (29.7,-160.1) {};
\node at (34,-193) {(b)};
\end{scope}
\end{tikzpicture}
\end{center}
\legende{Figure de l'édition, redessinée d'après le dessin de la page 43 (dessin 138/03-1), en deux panneaux. (a) La conique $C_{2}$ vue comme une sphère, triangulée en octaèdre par trois grands cercles : ses six sommets sont les deux points au-dessus de chacun des points $0$, $1$, $\infty$ de $C_{1}$, et l'axe qui joint les deux points au-dessus de $\infty$ est tracé ; le point sans étiquette sur l'équateur est sur la page. (b) La seconde figure, que la page ne commente pas et qui n'est lue qu'en partie : une courbe fermée portant $\infty$ en haut et en bas, un ovale intérieur portant deux points $0$, deux traits marqués $1$ et, en gris, un arc léger au crayon.}

\cote{138}{46}
\begin{center}
\begin{tikzpicture}[x=1cm,y=1cm,line cap=round,line join=round,
  pt/.style={circle,fill,inner sep=1.3pt},
  >={Stealth[length=1.8mm]}, font=\small]
% (a) les trois rangées et les flèches entre leurs points
\begin{scope}
\foreach \y in {0,1.4,2.8}
  \draw[thick] plot[smooth, tension=0.7] coordinates {(0,\y) (2,\y+0.08) (4.5,\y-0.06) (7,\y)};
% en bas
\node[pt,label={below:$\infty$}] (b-inf) at (1.6,0.03) {};
\node[pt,label={below:$0$}] (b-0) at (4.0,-0.04) {};
\node[pt,label={below:$1$}] (b-1) at (5.9,-0.03) {};
% au milieu
\node[pt,label={[inner sep=1pt]below left:$0\,(1)$}] (m-0) at (0.8,1.45) {};
\node[pt,label={[inner sep=1pt]below right:$\infty\,(1)$}] (m-inf) at (2.2,1.48) {};
\node[pt,label={[inner sep=1pt]below right:$1\,(2)$}] (m-1) at (4.0,1.36) {};
\node[pt,label={[inner sep=1pt]below right:$-1\,(2)$}] (m-m1) at (5.9,1.37) {};
% en haut
\node[pt,label={above:$0\,(n)$}] (h-0) at (0.8,2.85) {};
\node[pt,label={above:$\infty\,(n)$}] (h-inf) at (2.2,2.88) {};
\node[pt,label={[inner sep=1pt]above right:$1\,(1)$}] (h-1) at (4.0,2.76) {};
\node[pt,inner sep=0.9pt] at (4.0,3.15) {};
% les flèches
\draw[->,shorten <=2pt,shorten >=2pt] (h-0) -- (m-0);
\draw[->,shorten <=2pt,shorten >=2pt] (h-inf) -- (m-inf);
\draw[->,shorten <=2pt,shorten >=2pt] (m-0) .. controls (0.9,0.7) .. (b-inf);
\draw[->,shorten <=2pt,shorten >=2pt] (m-inf) .. controls (2.1,0.8) .. (b-inf);
\draw[->,shorten <=2pt,shorten >=2pt] (h-1) -- node[draw,rounded corners=5pt,fill=white,inner sep=1.5pt,font=\scriptsize] {$n$ fois} (m-1);
\draw[->,shorten <=2pt,shorten >=2pt] (m-1) -- (b-0);
\draw[->,shorten <=2pt,shorten >=2pt] (m-m1) -- (b-1);
\node at (7.6,2.1) {$n$ pair};
\node at (3.5,-0.85) {(b)};
\end{scope}
% (b) les deux lignes de points de la gauche de la page
\begin{scope}[shift={(-4.6,0.6)}]
\draw[thick] plot[smooth, tension=0.7] coordinates {(0,1.4) (1.6,1.46) (4.1,1.4)};
\draw[thick] plot[smooth, tension=0.7] coordinates {(0,0) (1.6,0.06) (4.1,0)};
\node[pt,label={above:$0\,(2)$}] at (0.4,1.42) {};
\node[pt,label={above:$1\,(2)$}] at (1.5,1.46) {};
\node[pt,label={[font=\scriptsize]above:$\infty_{1}\,(1)$}] at (2.45,1.44) {};
\node[pt,label={[font=\scriptsize]above:$\infty_{2}\,(1)$}] at (3.6,1.41) {};
\node[pt,label={below:$0$}] at (0.4,0.02) {};
\node[pt,label={below:$1$}] at (1.5,0.06) {};
\node[pt,label={below:$\infty$}] at (2.9,0.03) {};
\node at (1.8,-1.45) {(a)};
\end{scope}
\end{tikzpicture}
\end{center}
\legende{Figure de l'édition, redessinée d'après le dessin de la page 46 (dessin 138/03-3), en deux panneaux. (a) Les deux lignes de points de la gauche de la page : la fibre de $X_{1} \to X_{0}$ au-dessus de $0$, $1$, $\infty$. (b) Les fibres de la tour $X_{n} \to X_{1} \to X_{0}$, les trois rangées étant, de haut en bas, $X_{n}$, $X_{1}$ et $X_{0}$, l'indice de ramification entre parenthèses : au-dessus de $\infty$ les deux points $z = 0$ et $z = \infty$ de $X_{1}$, non ramifiés, eux-mêmes images de points ramifiés d'indice $n$ ; au-dessus de $0$ et de $1$ les points $z = 1$ et $z = -1$, ramifiés d'indice $2$, et au-dessus de $z = 1$ des points non ramifiés de $X_{n}$ (« $n$ fois »). Les pointes des deux flèches qui aboutissent à $\infty$ sont une lecture, et « $-1\,(2)$ » est en partie biffé sur la page.}

\cote{138}{48}
\begin{center}
\begin{tikzpicture}[x=1cm,y=1cm,line cap=round,line join=round,
  >={Stealth[length=1.8mm]}, font=\small]
\node at (0,0.75) {$X$};
\node at (3.6,0.75) {$X'$};
\node (x0) at (0,0) {$0$};
\node (x1) at (-0.35,-0.8) {$1$};
\node (xm1) at (0.3,-0.8) {$-1$};
\node at (0.95,-0.62) {$2$};
\node (xinf) at (0,-1.6) {$\infty$};
\node (y0) at (3.6,0) {$0$};
\node (y1) at (3.6,-0.8) {$1$};
\node (yinf) at (3.6,-1.6) {$\infty$};
\draw[->] (x0.east) .. controls (1.4,0.3) and (2.3,0.3) .. (y0.west);
\draw[->] (-0.05,-1.08) .. controls (1.2,-1.3) and (2.4,-1.05) .. (y1.west);
\draw[->] (xinf.east) .. controls (1.4,-1.9) and (2.3,-1.9) .. (yinf.west);
\draw[thick] (0,-2.1) -- (0,-3.3);
\draw[thick,->] (3.6,-2.1) -- node[right] {$n$} (3.6,-3.3);
\end{tikzpicture}
\end{center}
\legende{Figure de l'édition, redessinée d'après le dessin de la page 48 (dessin 138/03-5). Deux copies $X$ et $X'$ de la droite, et une application de $X$ dans $X'$ qui envoie $0$ sur $0$, $1$ et $-1$ sur $1$, $\infty$ sur $\infty$, dessinée à côté du calcul de valeurs critiques de la page ; au-dessous, les deux droites en traits verticaux, la seconde marquée $n$. Le $2$ écrit à droite de $1$, $-1$ est laissé à sa place, sans flèche.}


\section{Diagrammes commutatifs dessinés en perspective}

Grothendieck dessine souvent ses diagrammes en perspective : un carré de carrés devient un cube, une suite de cubes un prisme. Les figures qui suivent les redressent : les sommets aux coins d'un cube vu de biais, les arêtes de devant passant au-dessus de celles de derrière avec un blanc. Elles viennent des notes sur la spécialisation (cote 34), le produit des algèbres quadratiques (cote 84, où le double cube de la page 125 est le produit poussé et le produit tiré d'extensions $0 \to L \to E \to M \to 0$ qu'une section de l'article Lean énonce en coordonnées), les configurations de l'hexagone (cote 74), les sous-groupes du groupe de Teichmüller (cote 144), \emph{Vers une géométrie des formes} (cotes 156-4, 156-6, 156-9) et plusieurs autres liasses.

\cote{34}{143}
\begin{center}
\begin{tikzpicture}[x=1cm,y=1cm,line cap=round,>={Stealth[length=2mm]},
  every node/.style={inner sep=2pt},
  over/.style={preaction={draw=white,line width=3.2pt,arrows=-}}]
% front face
\node (V)  at (0,2.4) {$V$};
\node (Vb) at (3,2.4) {$\bar V$};
\node (S)  at (0,0)   {$S$};
\node (Sb) at (3,0)   {$\bar S$};
% back face
\node (Ve)  at (1.5,1.4)  {$V_\eta$};
\node (Veb) at (4.5,1.4)  {$V_{\bar\eta}$};
\node (e)   at (1.5,-1.0) {$\eta$};
\node (eb)  at (4.5,-1.0) {$\bar\eta$};
% the closed points
\node (s)  at (-0.9,0.75) {$s$};
\node (sb) at (2.25,0.8)  {$\bar s$};
% back edges first
\draw (V) -- (Ve);
\draw (Vb) -- (Veb);
\draw (Ve) -- (Veb);
\draw (Ve) -- (e);
\draw (Veb) -- (eb);
\draw[->] (eb) -- (e);
\draw[->] (e) -- (S);
\draw[->] (eb) -- (Sb);
% front edges, over the back ones
\draw[->] (Vb) -- (V);
\draw (V) -- (S);
\draw[->,over] (Vb) -- (Sb);
\draw[->,over] (Sb) -- (S);
\draw[->] (s) -- (S);
\draw[->] (sb) -- (Sb);
% the second V over S, at the right
\node (V2) at (6.2,1.6)  {$V$};
\node (S2) at (6.2,-0.2) {$S$};
\draw (V2) -- (S2);
\end{tikzpicture}
\end{center}
\legende{Figure de l'édition, redessinée d'après le dessin de la page 143 (dessin 34/08-1). Un cube de schémas au-dessus de $S$ : sur la face avant, $\bar V \to V$ au-dessus de $\bar S \to S$ ; sur la face arrière, $V_\eta$ et $V_{\bar\eta}$ au-dessus de $\eta \leftarrow \bar\eta$ ; les points fermés $s$ et $\bar s$ s'envoient dans $S$ et $\bar S$. Les traits sans pointe sont sans pointe sur la page ; à droite, un second $V$ au-dessus de $S$.}

\cote{74}{76}
\begin{center}
\begin{tikzpicture}[x=1cm,y=1cm,line cap=round,line join=round,
  >={Stealth[length=1.8mm]}, font=\footnotesize,
  ob/.style={align=center,inner sep=2pt},
  deg/.style={font=\scriptsize,inner sep=1.5pt},
  gap/.style={preaction={draw,white,line width=3pt,-}}]
% face du haut
\node[ob] (tau) at (0,0) {$\tau$ = diagonales};
\node[ob] (rep) at (5.5,0) {rep.\ sym.};
\node[ob] (som) at (3.2,2.6) {sommets};
\node[ob] (F) at (8.7,2.6) {sommets orientés $\simeq$\\ côtés $\simeq$ repères};
\node[ob] (cot) at (4.35,1.3) {côtés};
% face du bas
\node[ob] (e) at (0,-5) {$e$};
\node[ob] (om) at (5.5,-5) {$\omega = \omega(\tau)$\\ orientations};
\node[ob] (TI) at (2.85,-2.4) {triangles inscrits};
\node[ob] (eo) at (8.7,-2.4) {$\varepsilon \times \omega$\\ triangles inscrits orientés};
\node[ob] (ep) at (4.1,-3.75) {$\varepsilon' = \varepsilon \wedge \omega$\\ triangles coinscrits};
\node[ob] (eps) at (1.8,-1.4) {$\varepsilon$};
% arêtes arrière et intérieures
\draw[->] (eo) -- node[deg,above,pos=0.3] {$2$} (TI);
\draw[->] (F) -- node[deg,above] {$2$} (som);
\draw[->] (F) -- node[deg,above,pos=0.45] {$2$} (cot);
\draw[->] (F) -- (rep);
\draw[->] (F) -- (eo);
\draw[->] (eo) -- node[deg,above,pos=0.4] {$2$} (ep);
\draw[->] (eo) -- node[deg,below right] {$2$} (om);
\draw[->] (som) -- node[deg,above left] {$2$} (tau);
\draw[->] (TI) -- node[deg,above left] {$2$} (e);
\draw[->] (cot) -- node[deg,below right,pos=0.25] {$2$} (eps);
\draw[->] (eps) -- node[deg,left] {$2$} (e);
\draw[->,gap] (cot) -- (ep);
% arêtes avant
\draw[gap] (rep) -- (om);
\draw[->,gap] (rep) -- node[deg,below,pos=0.3] {$2$} (tau);
\draw[->] (tau) -- node[deg,left] {$3$} (e);
\draw[->] (om) -- node[deg,below] {$2$} (e);
\end{tikzpicture}
\end{center}
\legende{Figure de l'édition, redessinée d'après le dessin de la page 76 (dessin 74/04-9). Le « cube » des revêtements entre les ensembles attachés à un hexagone : en haut, les repères (« sommets orientés $\simeq$ côtés $\simeq$ repères ») et leurs images, sommets, côtés, « rep.\ sym. » et diagonales $\tau$ ; en bas, $\varepsilon \times \omega$, ses quotients $\varepsilon'$, triangles inscrits, orientations $\omega$, et le point $e$ ; les flèches portent le degré écrit sur la page, quand il l'est. Le coin supérieur gauche est surchargé : la répartition des degrés autour de $\tau$, l'origine des flèches qui partent de l'amas d'étiquettes en haut à droite, et celle de la flèche verticale qui descend vers $\varepsilon'$ sont des lectures ; la page note à part $\varepsilon$ et « triangles inscrits », et la figure les garde séparés.}

\cote{74}{79}
\begin{center}
\begin{tikzpicture}[x=1.15cm,y=1cm,line cap=round,line join=round,
  >={Stealth[length=1.8mm]}, font=\small,
  ob/.style={inner sep=2pt},
  deg/.style={font=\scriptsize,inner sep=1.5pt},
  gap/.style={preaction={draw,white,line width=3pt,-}}]
% le petit hexagone et ses deux triangles inscrits
\begin{scope}[shift={(-2.1,-2)},x=0.6cm,y=0.6cm]
\foreach \i in {0,...,5} \coordinate (h\i) at (60*\i:1);
\draw (h0) -- (h1) -- (h2) -- (h3) -- (h4) -- (h5) -- cycle;
\draw[densely dotted] (h0) -- (h2) -- (h4) -- cycle;
\draw[densely dotted] (h1) -- (h3) -- (h5) -- cycle;
\end{scope}
% étage du haut
\node[ob,label={[font=\scriptsize]right:$D_{6}$}] (Hr) at (7.5,0) {$H_{r}$};
\node[ob] (Hs) at (5.05,1.0) {$H_{s}$};
\node[ob,label={[font=\scriptsize,inner sep=1pt]above:$\mathfrak{S}_{3}$}] (Hp) at (3.75,0) {$H_{p'}$};
\node[ob] (Hss) at (2.45,-1.0) {$H_{s^{*}}$};
\node[ob] (Hd) at (0,0) {$H_{d}$};
% étage du bas
\node[ob,label={[font=\scriptsize]right:$\mathbb{F}_{2}^{2}$}] (Htr) at (7.5,-4) {$H_{tr}$};
\node[ob] (Htrr) at (5.05,-3) {$H_{trr}$};
\node[ob] (Hw) at (3.75,-4) {$H_{\omega}$};
\node[ob] (Htrs) at (2.45,-5) {$H_{tr^{*}}$};
\node[ob] (e) at (0,-4) {$e$};
% arrière
\draw[->] (Hr) -- node[deg,above] {$2$} (Hs);
\draw[->] (Htr) -- node[deg,above] {$2$} (Htrr);
\draw[->] (Hs) -- node[deg,right,pos=0.3] {$3$} (Htrr);
\draw[->] (Htrr) -- node[deg,above,pos=0.3] {$2$} (e);
% milieu
\draw[->,gap] (Hr) -- node[deg,above,pos=0.3] {$2$} (Hp);
\draw[->] (Htr) -- (Hw);
\draw[->,gap] (Hp) -- (Hw);
\draw[->] (Hp) -- node[deg,above] {$2$} (Hd);
\draw[->] (Hw) -- node[deg,above,pos=0.35] {$2$} (e);
% avant
\draw[->,gap] (Hr) -- node[deg,below,pos=0.2] {$2$} (Hss);
\draw[->] (Htr) -- (Htrs);
\draw[->,gap] (Hss) -- node[deg,right,pos=0.55] {$3$} (Htrs);
\draw[->] (Hss) -- node[deg,above right,pos=0.5] {$2$} node[font=\scriptsize,below left,pos=0.55] {opp} (Hd);
\draw[->] (Htrs) -- node[deg,below left] {$2$} (e);
\draw[->] (Hr) -- node[deg,right] {$3$} (Htr);
\draw[->] (Hd) -- node[deg,left] {$3$} (e);
\end{tikzpicture}
\end{center}
\legende{Figure de l'édition, redessinée d'après le dessin de la page 79 (dessin 74/04-12). Le cube des revêtements de l'hexagone : de $H_{r}$, marqué $D_{6}$, aux espaces homogènes $H_{s}$, $H_{p'}$, $H_{s^{*}}$, $H_{d}$ en haut et, en bas, $H_{tr}$ (marqué $\mathbb{F}_{2}^{2}$), $H_{trr}$, $H_{\omega}$, $H_{tr^{*}}$ et le point $e$ ; seuls les degrés écrits sur la page sont reportés. Les indices du milieu ($H_{p'}$, $H_{s^{*}}$, $H_{trr}$) sont surchargés et lus comme dans la transcription, et l'arrivée de deux des flèches de la rangée du haut en $H_{p'}$ ou en $H_{d}$ est en partie une lecture. À gauche, le petit hexagone de la page, ses deux triangles inscrits en pointillés.}

\cote{74}{90}
\begin{center}
\begin{tikzpicture}[x=1.05mm,y=1mm,line cap=round,line join=round,
  >={Stealth[length=1.6mm]}, font=\footnotesize,
  ob/.style={align=center,inner sep=1.5pt},
  sm/.style={font=\scriptsize},
  grp/.style={font=\scriptsize,inner sep=1pt},
  deg/.style={font=\scriptsize,inner sep=1.2pt},
  gap/.style={preaction={draw,white,line width=3pt,-}}]
\node[ob] (N1) at (58,-6) {$\mathrm{Car}_{d*} \simeq$ pav à trig pt};
\node[ob] (N2) at (28,-24) {Car\\ {\scriptsize $\simeq$ pav pt}};
\node[ob,label={[grp]right:$\mathfrak{S}_{3}$}] (N3) at (90,-24) {$\mathrm{Car}_{s*} \simeq$ pav.\ car\\ {\scriptsize $\simeq$ pav $+$ trig ép}};
\node[ob] (N5) at (10,-44) {$\mathrm{hex}_{\Omega}$\\ {\scriptsize $\simeq$ tria $\Omega$}};
\node[ob,label={[grp]above right:$\mathfrak{S}_{3}$}] (N4) at (60,-44) {$\mathrm{hex}_{r'} \simeq \mathrm{car}_{r'}$\\ {\scriptsize $\simeq$ tria ép}};
\node[ob] (N6) at (10,-64) {hex\\ {\scriptsize $\simeq$ tria}};
\node[ob,label={[grp]right:$\mathfrak{S}_{3}$}] (N7) at (28,-72) {$\mathrm{hex}_{\mathrm{tr}*}$};
\node[ob] (N8) at (48,-64) {$\mathrm{hex}_{\mathrm{tr}\Omega}$\\ {\scriptsize ($\simeq$ trig ép tria $\Omega$)}};
\node[ob,label={[grp]right:$\mathfrak{S}_{2}$}] (N9) at (92,-64) {hex ép\\ {\scriptsize ($\simeq$ trig ép tria ép)}\\ {\scriptsize ($\simeq$ pav ép)}};
\node[ob] (N10) at (66,-80) {$\mathrm{hex}_{d}$\\ {\scriptsize ($\simeq \mathrm{Car}_{d} \simeq$ tria pt)}};
\node[ob] (N11) at (92,-80) {$\mathrm{hex}_{s*}$};
\node[ob] (N12) at (10,-96) {Pav coor\\ {\scriptsize $\simeq$ trig tria}};
\node[ob,label={[grp]below left:$\mathfrak{S}_{3}$}] (N14) at (40,-96) {$\mathrm{hex}_{\mathrm{tr}}$\\ {\scriptsize ($\simeq$ trig ép tria)}};
\node[ob,label={[grp]right:$\mathfrak{S}_{3}$}] (N15) at (92,-96) {$\mathrm{hex}_{s}$\\ {\scriptsize ($\simeq$ tria ép)}};
\node[ob] (N16) at (66,-110) {trig tria pt\\ {\scriptsize $\simeq$ tria}};
\node[ob] (N13) at (10,-124) {Pav};
\node[ob] (N17) at (40,-124) {trig $\Omega$ tria};
\node[ob] (N18) at (92,-124) {tria $\Omega$};
% flèches
\draw[->] (N2) -- node[deg,above left] {$2$} (N1);
\draw[->] (N3) -- node[deg,above right] {$2$} (N1);
\draw[->] (N9) -- (N3);
\draw[->] (N4) -- node[deg,above] {$3$} (N5);
\draw[->] (N9) -- node[deg,below right] {$2$} (N4);
\draw[->] (N5) -- node[deg,left] {$2$} (N6);
\draw[->] (N5) -- node[deg,above right,pos=0.45] {$2$} (N8);
\draw[->] (N9) -- node[deg,above] {$3$} (N8);
\draw[->] (N8) -- (N7);
\draw[->] (N7) -- (N6);
\draw[->] (N6) -- node[deg,left] {$3$} (N12);
\draw[->] (N6) -- node[deg,above right,pos=0.6] {$2$} (N14);
\draw[->] (N8) -- node[deg,right,pos=0.6] {$2$} (N14);
\draw[->] (N11) -- (N10);
\draw[->] (N15) -- node[deg,below left] {$3$} (N10);
\draw[->] (N15) -- node[deg,above,pos=0.3] {$3$} (N14);
\draw[->,gap] (N10) -- node[deg,right,pos=0.25] {$3$} (N16);
\draw[->] (N14) -- node[deg,above] {$3$} (N12);
\draw[->] (N12) -- node[deg,left] {$2$} (N13);
\draw[->] (N12) -- node[deg,above right] {$2$} (N17);
\draw[->] (N15) -- (N18);
\draw[->] (N16) -- node[deg,above right] {$2$} (N18);
\draw[->] (N18) -- node[deg,above] {$3$} (N17);
\end{tikzpicture}
\end{center}
\legende{Figure de l'édition, redessinée d'après le dessin de la page 90 (dessin 74/05-16). Le diagramme des ensembles de carrés et d'hexagones d'un pavé, avec leurs groupes : chaque flèche porte le degré écrit sur la page, quand il l'est. Les ensembles sont replacés sur une grille, sans changer les flèches du dessin ; les étiquettes sont celles de la transcription, « trig pt » (en haut) et « car » (dans $\mathrm{Car}_{s*}$) étant des lectures douteuses. Le dessin est au crayon léger : le sens de plusieurs diagonales, et l'origine de la flèche de degré $2$ qui monte vers $\mathrm{Car}_{d*}$, sont des lectures.}

\cote{84}{125}
\begin{center}
\begin{tikzpicture}[x=1.35mm,y=1.35mm,line cap=round,>={Stealth[length=1.8mm]},
  every node/.style={inner sep=2pt},
  eq/.style={double,double distance=1.3pt},
  gap/.style={white,line width=4.5pt}]
% back face (shifted by the depth vector (10,12))
\node (L)   at (10,34) {$L_i$};
\node (E)   at (38,34) {$E_i'$};
\node (Mp)  at (66,34) {$M_i'$};
\node (L3)  at (10,12) {$L_i'''$};
\node (E3)  at (38,12) {$E_i'''$};
\node (Mp3) at (66,12) {$M_i'$};
% front face
\node (L2)  at (0,22)  {$L_i''$};
\node (E2)  at (28,22) {$E_i''$};
\node (M)   at (56,22) {$M_i$};
\node (L1)  at (0,0)   {$L_i'$};
\node (E1)  at (28,0)  {$E_i'$};
\node (M1)  at (56,0)  {$M_i$};
% back face edges
\draw (L) -- (E);
\draw[->] (E) -- (Mp);
\draw (L3) -- (E3);
\draw[->] (Mp3) -- (E3);
\draw (L) -- (L3);
\draw[->] (E) -- (E3);
\draw[eq] (Mp) -- (Mp3);
% depth edges, back to front
\draw[eq] (L) -- (L2);
\draw[->] (E) -- (E2);
\draw[->] (Mp) -- (M);
\draw[->] (L3) -- (L1);
\draw[->] (E3) -- (E1);
\draw (Mp3) -- (M1);
% front face edges, over the back ones
\draw[gap] (L2) -- (E2);
\draw (L2) -- (E2);
\draw[gap] (E2) -- (M);
\draw (E2) -- (M);
\draw[->] (L1) -- (E1);
\draw[->] (E1) -- (M1);
\draw[eq,-{Implies}] (L2) -- (L1);
\draw[gap] (E2) -- (E1);
\draw[->] (E2) -- (E1);
\draw[gap] (M) -- (M1);
\draw[eq] (M) -- (M1);
\end{tikzpicture}
\end{center}
\legende{Figure de l'édition, redessinée d'après le dessin de la page 125 (dessin 84/07-1). Deux cubes partageant la face des $E$ : les suites $L_i \to E_i \to M_i$ et leurs copies primées, placées aux sommets, sur trois colonnes $L$, $E$, $M$ ; face avant $L_i''$, $E_i''$, $M_i$ au-dessus de $L_i'$, $E_i'$, $M_i$, face arrière $L_i$, $E_i'$, $M_i'$ au-dessus de $L_i'''$, $E_i'''$, $M_i'$. Les traits doubles, les traits sans pointe et les primes des étiquettes sont ceux du dessin ; les deux traits doubles verticaux de droite sont les signes d'égalité de la page, et plusieurs primes restent une lecture.}

\cote{85}{97}
\begin{center}
\begin{tikzpicture}[x=2.1mm,y=2.1mm,line cap=round,line join=round,>={Stealth[length=1.6mm,width=1.2mm]},font=\small,every node/.style={inner sep=2pt},shorten >=0.6pt,shorten <=0.6pt,]
\node[circle,fill,inner sep=0pt,minimum size=2.6pt] (p1) at (12.3,-33.6) {};
\node[circle,fill,inner sep=0pt,minimum size=2.6pt] (p2) at (12.4,-24.6) {};
\node[circle,fill,inner sep=0pt,minimum size=2.6pt] (p3) at (17.1,-26.4) {};
\node[circle,fill,inner sep=0pt,minimum size=2.6pt] (p4) at (17.3,-17.1) {};
\node[circle,fill,inner sep=0pt,minimum size=2.6pt] (p5) at (19.6,-18) {};
\node[circle,fill,inner sep=0pt,minimum size=2.6pt] (p6) at (22.1,-18.8) {};
\node[circle,fill,inner sep=0pt,minimum size=2.6pt] (p7) at (22.3,-9.8) {};
\node[circle,fill,inner sep=0pt,minimum size=2.6pt] (p8) at (28.8,-24.5) {};
\node[circle,fill,inner sep=0pt,minimum size=2.6pt] (p9) at (28.4,-33.5) {};
\node[circle,fill,inner sep=0pt,minimum size=2.6pt] (p10) at (32.5,-26.6) {};
\node[circle,fill,inner sep=0pt,minimum size=2.6pt] (p11) at (34.7,-26.4) {};
\node[circle,fill,inner sep=0pt,minimum size=2.6pt] (p12) at (34.6,-17.4) {};
\node[circle,fill,inner sep=0pt,minimum size=2.6pt] (p13) at (38.4,-18.5) {};
\node[circle,fill,inner sep=0pt,minimum size=2.6pt] (p14) at (38.5,-9.8) {};
\node[circle,fill,inner sep=0pt,minimum size=2.6pt] (p15) at (46.6,-30.3) {};
\node[circle,fill,inner sep=0pt,minimum size=2.6pt] (p16) at (45.6,-39.5) {};
\node[circle,fill,inner sep=0pt,minimum size=2.6pt] (p17) at (25,-47.5) {};
\node[circle,fill,inner sep=0pt,minimum size=2.6pt] (p18) at (20.6,-55.1) {};
\node[inner sep=0pt,minimum size=0pt] (T1) at (41.1,-36.2) {};
\node[inner sep=0pt,minimum size=0pt] (T2) at (41,-45.8) {};
\node[inner sep=0pt,minimum size=0pt] (Td) at (36.4,-54.9) {};
\node[inner sep=0pt,minimum size=0pt] (Bd) at (33.1,-28.2) {};
\node[inner sep=0pt,minimum size=0pt] (Cv) at (32.5,-39.9) {};
\node[inner sep=0pt,minimum size=0pt] (Pd) at (29.1,-40.1) {};
\draw[->] (p1) -- (p2);
\draw[->] (p1) -- (p3);
\draw[->] (p1) -- (p18);
\draw[->] (p2) -- (p8);
\draw[->] (p2) -- (p4);
\draw[->] (p3) -- (p4);
\draw[->] (p3) -- (p5);
\draw[->] (p3) -- (p6);
\draw[->] (p3) -- (p10);
\draw[->] (p3) -- (p17);
\draw[->] (p5) -- (p7);
\draw[->] (p6) -- (p7);
\draw[->] (p4) -- (p7);
\draw[->] (p7) -- (p14);
\draw[->] (p6) -- (p13);
\draw[->] (p9) -- (p8);
\draw[->] (p9) -- (p10);
\draw[->] (p9) -- (p11);
\draw[->] (p9) -- (p13);
\draw[->] (p12) -- (p14);
\draw[->] (p13) -- (p14);
\draw[->] (p14) -- (p15);
\draw[->] (p16) -- (p15);
\draw[->] (p18) -- (p17);
\draw[-] (p4) -- (p12);
\draw[-] (p13) -- (T1);
\draw[->] (T1) -- (p15);
\draw[-] (p11) -- (T2);
\draw[-] (Bd) -- (T2);
\draw[->] (T2) -- (p16);
\draw[-] (p17) .. controls (27.2,-43) and (29.5,-40.6) .. (Cv);
\draw[->] (Cv) -- (p16);
\draw[-,densely dashed] (p18) -- (Td);
\draw[-,densely dashed] (p9) -- (Td);
\draw[-,densely dashed] (T2) -- (Td);
\draw[-,densely dashed,gray] (p6) -- (Pd);
\draw[->,preaction={draw,white,line width=2.5pt,-}] (p1) -- (p9);
\draw[->,preaction={draw,white,line width=2.5pt,-}] (p8) -- (p12);
\draw[->,preaction={draw,white,line width=2.5pt,-}] (p11) -- (p12);
\draw[-,preaction={draw,white,line width=2.5pt,-}] (p12) -- (T1);
\node[font=\small,inner sep=1pt] at (22.3,-7.9) {$\mathbb{D}_4$};
\node[font=\small,inner sep=1pt] at (39.7,-8) {$\mathfrak{S}_4$};
\node[font=\small,inner sep=1pt] at (19.9,-16.2) {$\mathbb{Z}_4$};
\node[font=\small,inner sep=1pt] at (38,-21.6) {$\mathfrak{A}_4$};
\node[font=\small,inner sep=1pt] at (10.3,-23.6) {$\mathfrak{S}_2$};
\node[font=\small,inner sep=1pt] at (27.8,-22.6) {$\mathfrak{S}_3$};
\node[font=\small,inner sep=1pt] at (18.9,-28.1) {$\mathbb{F}_2$};
\node[font=\small,inner sep=1pt] at (32.2,-24.8) {$\mathbb{Z}_6$};
\node[font=\small,inner sep=1pt] at (36.6,-26.6) {$\mathbb{D}_3$};
\node[font=\small,inner sep=1pt] at (8.4,-34.2) {$\mathfrak{A}_2=1$};
\node[font=\small,inner sep=1pt] at (31.6,-35.3) {$\mathfrak{A}_3=\mathbb{Z}_3$};
\node[font=\small,inner sep=1pt] at (48.8,-29.8) {$\mathfrak{S}_5$};
\node[font=\small,inner sep=1pt] at (48.2,-40.5) {$\mathfrak{A}_5$};
\node[font=\small,inner sep=1pt] at (27.4,-48.9) {$\mathbb{D}_5$};
\node[font=\small,inner sep=1pt] at (18.5,-56.6) {$\mathbb{Z}_5$};
\node[font=\small,inner sep=1pt,anchor=east] at (16.6,-15.6) {$\mathfrak{S}_2\times\mathfrak{S}'_2=\mathfrak{S}_2\times\mathbb{F}_2$};
\node[font=\small,inner sep=1pt,anchor=west] at (22.9,-20.1) {$\mathbb{D}_2=\mathbb{F}_2^{2}$};
\node[font=\small,inner sep=1pt,anchor=west] at (35.4,-15.3) {$\mathfrak{S}_3\times\mathfrak{S}''_2$};
\node[font=\scriptsize,inner sep=1pt] at (11.1,-28.9) {$2$};
\node[font=\scriptsize,inner sep=1pt] at (14.6,-29.4) {$2$};
\node[font=\scriptsize,inner sep=1pt] at (20.9,-22.4) {$2$};
\node[font=\scriptsize,inner sep=1pt] at (22.9,-32.7) {$3$};
\node[font=\scriptsize,inner sep=1pt] at (14.8,-44.8) {$5$};
\end{tikzpicture}
\end{center}
\legende{Figure de l'édition, redessinée d'après le dessin de la page 97 (dessin 85/05-3). Le treillis en perspective des classes de sous-groupes de $\mathfrak{S}_5$, avec les indices $2$, $3$, $5$ qu'il écrit ; les flèches sont des inclusions. Trois flèches descendent à un point de rebroussement avant de remonter vers $\mathfrak{S}_5$ et $\mathfrak{A}_5$, comme sur la page, et les traits pointillés sont des arêtes cachées ; ce tracé en V est une lecture. Les correspondances notées sous le treillis ne sont pas reprises.}

\cote{117}{24}
\begin{center}
\begin{tikzpicture}[x=1mm,y=1mm,line cap=round,line join=round,>={Stealth[length=1.6mm,width=1.2mm]},font=\small,every node/.style={inner sep=2pt},]
\node[circle,draw,minimum size=5.5mm,inner sep=0pt] (R1) at (58.5,-16.3) {$\cdot$};
\node[circle,draw,minimum size=5.5mm,inner sep=0pt] (R2) at (37.5,-34.5) {$e$};
\node[circle,fill,inner sep=0.9pt] (J1) at (57,-34.5) {};
\node[circle,fill,inner sep=0.9pt] (J0) at (76.1,-32.4) {};
\node[circle,draw,minimum size=5.5mm,inner sep=0pt] (R3) at (14,-55) {$e$};
\node[circle,fill,inner sep=0.9pt] (J2) at (37.5,-55) {};
\node[circle,fill,inner sep=0.9pt] (J3) at (57,-55) {};
\node[circle,draw,minimum size=5.5mm,inner sep=0pt] (R4) at (78.5,-55) {$e$};
\node[circle,draw,minimum size=5.5mm,inner sep=0pt,label=left:$X$] (X) at (13.5,-74.5) {};
\node[circle,fill,inner sep=0.9pt] (J4) at (37.5,-74.5) {};
\node[circle,draw,minimum size=5.5mm,inner sep=0pt] (R6) at (58,-74.5) {$e$};
\node[circle,draw,minimum size=5.5mm,inner sep=0pt,label=left:$Y$] (Y) at (13,-92.4) {};
\node[circle,draw,minimum size=5.5mm,inner sep=0pt] (R8) at (36.5,-92.4) {$e$};
\draw[->] (J1) -- (R2);
\draw[->] (J0) -- (J1);
\draw[->] (J3) -- (J2);
\draw[->] (R4) -- (J3);
\draw[->] (J4) -- (X);
\draw[->] (R6) -- (J4);
\draw[->] (R8) -- (Y);
\draw[->,thick] (J2) -- (R3);
\draw[->] (R2) -- (J2);
\draw[->] (J1) -- (J3);
\draw[->] (J2) -- (J4);
\draw[->] (J3) -- (R6);
\draw[->] (R3) -- (X);
\draw[->] (X) -- (Y);
\draw[->] (J4) -- (R8);
\draw[->] (R1) -- (R2);
\draw[->] (R2) -- (R3);
\draw[->] (R4) -- (R6);
\draw[->] (R6) -- (R8);
\draw[line width=1.2pt] (R1) -- (J1);
\draw[line width=1.2pt] (60.5,-9) -- (61.5,-32.4) -- (J0) -- (R4);
\draw[line width=1.2pt] (79.5,-33) -- (80.5,-40) -- (82.5,-80);
\draw[gray] (76,0) -- (J0) -- (76.6,-90);
\draw[gray] (J4) -- (J1);
\draw (91,-37) -- (80.5,-37);
\draw (24,-10.7) -- (58,-9);
\draw[gray] (26.5,-23.5) -- (80,-21.5);
\draw (62,-19) -- (74,-18.7);
\draw (81,-16.5) -- (88.5,-17);
\draw (94,-16) -- (94,-31);
\end{tikzpicture}
\end{center}
\legende{Figure de l'édition, redessinée d'après le dessin de la page 24 (dessin 117/02-1). Le croquis en perspective de l'escalier : les sommets cerclés, dont cinq marqués $e$, l'objet nul, et $X$, $Y$ en bas à gauche ; les points noirs sont des sommets sans lettre, et les flèches vont vers la gauche, vers le bas et de l'arrière vers l'avant. Le trait épais en escalier, l'axe vertical et les traits du haut (en gris ceux qui sont au crayon) sont reportés tels quels, sans interprétation.}

\cote{119}{149}
\begin{center}
\begin{tikzpicture}[x=1.2mm,y=1.2mm,line cap=round,line join=round,>={Stealth[length=1.6mm,width=1.2mm]},font=\small,every node/.style={inner sep=2pt},]
\node (PBt) at (64,-8) {$\mathrm{PB}$};
\node (PBCR) at (24,-36) {$\mathrm{PBCR}$};
\node (EPCR) at (54,-33) {$\mathrm{EPCR}$};
\node (EPC) at (71,-44) {$\mathrm{EPC}$};
\node (PBC) at (40,-52) {$\mathrm{PBC}$};
\node (PBb) at (63,-66) {$\mathrm{PB}$};
\node (EP) at (88,-58) {$\mathrm{EP}$};
\draw[->] (PBt) -- node[pos=0.45,sloped,above,font=\scriptsize,inner sep=1.5pt] {$\text{recoller cône}$} (EPCR);
\draw[->] (PBt) -- node[pos=0.33,sloped,above,font=\scriptsize,inner sep=1.5pt] {$\text{identité}$} (PBb);
\draw[->,double,double distance=1.4pt,-{Implies}] (PBt) -- node[pos=0.5,sloped,above,font=\scriptsize,inner sep=1.5pt] {$\text{recoller collier}$} (PBCR);
\draw[->] (PBt) -- node[pos=0.5,sloped,above,font=\scriptsize,inner sep=1.5pt] {$\text{contraction }B$} (EP);
\draw[->] (PBCR) -- (EPCR);
\draw[->] (PBC) -- (PBb);
\draw[->] (PBb) -- (EP);
\draw[->] (PBCR) -- (PBC);
\draw[->] (EPC) -- (EP);
\draw[->,preaction={draw,white,line width=3pt,-}] (EPCR) -- (EPC);
\draw[->,preaction={draw,white,line width=3pt,-}] (PBC) -- (EPC);
\end{tikzpicture}
\end{center}
\legende{Figure de l'édition, redessinée d'après le dessin de la page 149 (dessin 119/08-1). Le système des catégories redessiné en perspective : de $\mathrm{PB}$ partent « recoller collier » (flèche double) vers $\mathrm{PBCR}$, « recoller cône » vers $\mathrm{EPCR}$, « identité » vers $\mathrm{PB}$, qui passe derrière $\mathrm{EPC}$, et « contraction $B$ » vers $\mathrm{EP}$. Le sommet du bas, qu'on lirait aussi « PD », est $\mathrm{PB}$ ; les petites marques illisibles entre $\mathrm{EPC}$ et $\mathrm{EP}$ ne sont pas reprises.}

\cote{122}{16}
\begin{center}
\begin{tikzpicture}[x=1cm,y=1cm,line cap=round,>={Stealth[length=1.8mm]},
  every node/.style={inner sep=2pt},
  hk/.style={{Hooks[right]}->},
  dot/.style={circle,fill,inner sep=0.8pt},
  ring/.style={circle,draw,fill=white,inner sep=1pt},
  oval/.style={gray,thin}]
% (a) the model
\node (Stop) at (1.7,4.8) {$S$};
\node (W)    at (4.3,4.8) {$W$};
\node (Ttop) at (6.9,4.8) {$T$};
\node (s)    at (0.4,3.7) {$s$};
\node (P)    at (3.0,3.7) {$P$};
\node (Q)    at (5.6,3.7) {$Q$};
\node (t)    at (8.2,3.7) {$t$};
\node (M)    at (1.7,2.6) {$M$};
\node (R)    at (4.3,2.8) {$R$};
\node[inner sep=0.5pt] at (4.3,2.43) {\rotatebox{90}{$=$}};
\node (TS)   at (4.3,2.07) {$T \times S$};
\node (N)    at (6.9,2.6) {$N$};
\node (T)    at (3.0,1.15) {$\boldsymbol{T}$};
\node (S)    at (5.6,1.15) {$\boldsymbol{S}$};
\node (st)   at (4.3,0.25) {$(s,t)$};
\node (e)    at (4.3,-0.75) {$\emptyset$};
\node (tl)   at (6.8,0.0) {$t$};
\node (Tl)   at (8.3,0.95) {$T$};
\node at (0.15,1.45) {$s$};
\node at (1.35,-0.05) {$s$};
% pencil ovals
\draw[oval,rotate around={-36:(3.65,3.25)}] (3.65,3.25) ellipse[x radius=1.2,y radius=0.42];
\draw[oval,rotate around={36:(4.95,3.25)}] (4.95,3.25) ellipse[x radius=1.2,y radius=0.42];
\draw[oval,rotate around={-47:(2.35,1.88)}] (2.35,1.88) ellipse[x radius=1.15,y radius=0.42];
\draw[oval,rotate around={47:(6.25,1.88)}] (6.25,1.88) ellipse[x radius=1.15,y radius=0.42];
% hatchings
\fill[pattern=north east lines,pattern color=gray] (1.05,0.75) ellipse[x radius=0.75,y radius=0.45];
\fill[pattern=north east lines,pattern color=gray,rotate around={35:(7.45,1.2)}] (7.45,1.2) ellipse[x radius=0.65,y radius=0.3];
% arrows
\draw[gray] (s) -- (Stop);
\draw[->] (M) -- ($(M)!0.78!(s)$);
\draw[->] (P) -- (Stop);
\draw[hk] (P) -- node[dot,pos=0.45] {} (W);
\draw[hk] (Q) -- node[dot,pos=0.45] {} ($(Q)!0.72!(W)$);
\draw[->] (Q) -- (Ttop);
\draw[->] (t) -- (Ttop);
\draw[hk] (M) -- node[dot] {} (P);
\draw[hk] (R) -- node[ring] {} (P);
\draw[hk] (R) -- node[dot] {} (Q);
\draw[hk] (N) -- (Q);
\draw[->] (N) -- (t);
\draw[hk] (T) -- node[font=\scriptsize,inner sep=0pt] {$\times$} (M);
\draw[hk] (T) -- node[dot] {} (TS);
\draw[hk] (S) -- node[dot] {} (TS);
\draw[hk] (S) -- node[font=\scriptsize,fill=white,inner sep=0.3pt] {$\ast$} (N);
\draw[hk] (st) -- (T);
\draw[hk] (st) -- (S);
\draw[hk] (e) -- (st);
\draw[->] (S) -- (tl);
\draw[dashed] (tl) -- (Tl);
\node at (4.3,-1.5) {(a)};
\end{tikzpicture}

\bigskip

\begin{tikzpicture}[x=1cm,y=1cm,line cap=round,>={Stealth[length=1.8mm]},
  every node/.style={inner sep=2pt},
  hk/.style={{Hooks[right]}->}]
% (b) its geometric counterpart
\coordinate (Stop) at (1.7,4.8);
\node (W)    at (4.3,4.8) {$V'$};
\node (Ttop) at (6.9,4.8) {$V'_s$};
\node (s)    at (0.4,3.7) {$s$};
\node (P)    at (3.0,3.7) {$X'_\eta$};
\node (Q)    at (5.6,3.7) {$V'_s$};
\node (t)    at (8.2,3.7) {$t$};
\node (M)    at (1.7,2.6) {$X'_\eta$};
\node (R)    at (4.3,2.6) {$V'_s \times \eta$};
\node (N)    at (6.9,2.6) {$e$};
\node (T)    at (3.0,1.15) {$V'_s$};
\node (S)    at (5.6,1.15) {$\eta$};
\node (st)   at (4.3,0.25) {$(s,t)$};
% frame round X'_eta and V'_s x eta, pencil frame round the lower X'_eta and V'_s
\draw[double,double distance=1.2pt,rounded corners=3pt,rotate around={-40:(3.65,3.15)}] (1.75,2.55) rectangle (5.55,3.75);
\draw[gray,thin,rounded corners=3pt,rotate around={-47:(2.35,1.88)}] (0.9,1.43) rectangle (3.8,2.33);
% hatched quadrilateral at the right, with the chevron at its corner and the loop beside it
\draw[pattern=north east lines,pattern color=gray] (6.05,0.55) -- (6.75,-0.25) -- (7.95,0.55) -- (7.25,1.35) -- cycle;
\node[fill=white,inner sep=1pt] (tl) at (6.95,0.35) {$t$};
\draw (8.05,0.8) -- (8.25,0.55) -- (8.0,0.4);
\draw (8.3,-0.05) circle[radius=0.1];
% arrows
\draw[->] (s) -- (Stop);
\draw[->] (P) -- (Stop);
\draw[->] (M) -- (s);
\draw[hk] (P) -- (W);
\draw[hk] (Q) -- (W);
\draw[->] (Q) -- node[pos=0.25,above left,inner sep=0.5pt] {$\scriptstyle\sim$} (Ttop);
\draw[->] (t) -- (Ttop);
\draw[->] (N) -- node[pos=0.5,above left,inner sep=0.5pt] {$\scriptstyle\sim$} (t);
\draw[hk] (N) -- (Q);
\draw[->] (M) -- (P);
\draw[->] (T) -- (M);
\draw[->] (T) -- (R);
\draw[->] (R) -- (P);
\draw[->] (R) -- (Q);
\draw[->] (S) -- (R);
\draw[->] (S) -- (N);
\draw[->] (st) -- (T);
\draw[->] (st) -- (S);
\node at (4.3,-0.75) {(b)};
\end{tikzpicture}
\end{center}
\legende{Figure de l'édition, redessinée d'après les deux dessins de la page 16 (dessins 122/01-5 et 01-6), réunis en deux panneaux sur une même disposition. (a) Le treillis des morceaux du modèle, de $\emptyset$ et du point-base $(s,t)$ jusqu'à $W$ en passant par $\boldsymbol{T}$, $\boldsymbol{S}$, $M$, $R = T \times S$, $N$, puis $P$, $Q$, et vers les bases $S$, $T$ ; les ovales gris sont au crayon sur la page, les points, l'anneau, la croix et l'astérisque sont portés par les flèches, et deux plages hachurées sont figurées en gris. Ce dessin est partiel : les flèches de $M$ vers $s$ et de $Q$ vers le haut s'arrêtent sans atteindre de lettre sûre. (b) Les objets correspondants selon le dictionnaire de la page : $V'$, $X'_\eta$ (deux fois), $V'_s$, $V'_s \times \eta$, $e$, $\eta$ ; le sommet en haut à gauche, raturé sur la page, est laissé sans nom, et une lettre biffée à droite n'est pas reprise.}

\cote{144}{133}
\begin{center}
\begin{tikzpicture}[x=1.05cm,y=1.05cm,line cap=round,line join=round,>={Stealth[length=1.6mm,width=1.2mm]},font=\small,every node/.style={inner sep=2pt},]
\node (mt) at (0,2.8) {$\mu_{\tau'}$};
\node (Dt) at (2.6,2.8) {$D_{\tau'}$};
\node (one) at (0,0) {$\{1\}$};
\node (mu) at (2.6,0) {$\underline{\mu}$};
\node (Ds) at (5.6,4.3) {$D_\sigma$};
\node (Pt) at (6.5,2.8) {$\Pi_{\tau'}$};
\node (Ls) at (5.8,1.95) {$L_\sigma$};
\node (Pi) at (6.5,0) {$\underline{\Pi}$};
\node (Dr) at (7.9,0.75) {$D_\rho$};
\node (Lr) at (7.9,-1.7) {$L_\rho$};
\node (S) at (10.2,2.8) {$\mathfrak{S}$};
\node (Sp) at (10.2,0) {$\underline{\mathfrak{S}}^+$};
\draw[{Hooks[right]}->] (Dt) -- node[midway,fill=white,font=\scriptsize,inner sep=0.8pt] {$2$} (Ds);
\draw[{Hooks[right]}->] (mu) -- node[midway,fill=white,font=\scriptsize,inner sep=0.8pt] {$2$} (Ls);
\draw[-] (Ls) -- node[midway,fill=white,font=\scriptsize,inner sep=0.8pt] {$2$} (Pt);
\draw[-] (Ls) -- node[midway,fill=white,font=\scriptsize,inner sep=0.8pt] {$2$} (Pi);
\draw[{Hooks[right]}->] (Ds) -- (S);
\draw[{Hooks[right]}->] (mt) -- node[midway,fill=white,font=\scriptsize,inner sep=0.8pt] {$2$} (Dt);
\draw[{Hooks[right]}->] (one) -- node[midway,fill=white,font=\scriptsize,inner sep=0.8pt] {$2$} (mu);
\draw[-] (one) -- node[pos=0.7,fill=white,font=\scriptsize,inner sep=0.8pt] {$2$} (mt);
\draw[-] (mu) -- node[midway,fill=white,font=\scriptsize,inner sep=0.8pt] {$2$} (Dt);
\draw[{Hooks[right]}->] (Dt) -- (Pt);
\draw[{Hooks[right]}->] (mu) -- (Pi);
\draw[{Hooks[right]}->] (Pt) -- node[midway,fill=white,font=\scriptsize,inner sep=0.8pt] {$6$} (S);
\draw[{Hooks[right]}->] (Pi) -- node[midway,fill=white,font=\scriptsize,inner sep=0.8pt] {$6$} (Sp);
\draw[-] (S) -- node[midway,fill=white,font=\scriptsize,inner sep=0.8pt] {$2$} (Sp);
\draw[{Hooks[right]}->] (mu) -- node[midway,fill=white,font=\scriptsize,inner sep=0.8pt] {$3$} (Lr);
\draw[{Hooks[right]}->] (Dr) -- (S);
\draw[{Hooks[right]}->] (Dr) -- node[midway,fill=white,font=\scriptsize,inner sep=0.8pt] {$6$} (Sp);
\draw[{Hooks[right]}->] (Lr) -- (Sp);
\draw[{Hooks[right]}->,preaction={draw,white,line width=3pt,-}] (Dt) -- node[pos=0.3,fill=white,font=\scriptsize,inner sep=0.8pt] {$3$} (Dr);
\draw[-,preaction={draw,white,line width=3pt,-}] (Lr) -- node[midway,fill=white,font=\scriptsize,inner sep=0.8pt] {$2$} (Dr);
\end{tikzpicture}
\end{center}
\legende{Figure de l'édition, redessinée d'après le dessin de la page 133 (dessin 144/07-4). Le diagramme (I), position relative de $L_\rho$ et $L_\sigma$ : les inclusions entre $\{1\}$, $\mu$, $\mu_{\tau'}$, $D_{\tau'}$, les sous-groupes $L_\sigma$, $L_\rho$, leurs normalisateurs diédraux $D_\sigma$, $D_\rho$, et $\Pi_{\tau'}$, $\underline{\Pi}$, $\mathfrak{S}$, $\underline{\mathfrak{S}}^+$, chaque flèche à crochet portant son indice ; les traits sans flèche joignent les groupes que la page relie d'un simple trait. Les étiquettes au crayon $SD_\sigma$, $SZ_\sigma$, trop pâles pour être placées, ne sont pas reportées.}

\cote{144}{133}
\begin{center}
\begin{tikzpicture}[x=1cm,y=1cm,line cap=round,line join=round,>={Stealth[length=1.6mm,width=1.2mm]},font=\small,every node/.style={inner sep=2pt},]
\node (D2) at (0,0) {$\mu\times\mu_\tau=D_2$};
\node[label={[font=\scriptsize,inner sep=1pt]above:$\mu\times\Pi'_0=$}] (Sh) at (4.2,0) {$\mathfrak{S}^!=\Pi^!_0$};
\node (S) at (7.6,0) {$\mathfrak{S}$};
\node (mu) at (0,-3) {$\underline{\mu}$};
\node[label={[font=\scriptsize,inner sep=1pt,xshift=-0.55cm]below:$\mu\times\Pi_0=$}] (Sp) at (4.2,-3) {$\mathfrak{S}^{!+}=\Pi$};
\node (Spp) at (7.6,-3) {$\underline{\mathfrak{S}}^+$};
\node (mt) at (1.7,-1.15) {$\mu_\tau$};
\node (Sh0) at (5.9,-1.15) {$\mathfrak{S}^!_0=\Pi'$};
\node (one) at (1.7,-4.15) {$\{1\}$};
\node (P0) at (5.9,-4.15) {$\mathfrak{S}^{!+}_0=\Pi_0$};
\draw[{Hooks[right]}->] (mt) -- node[pos=0.35,fill=white,font=\scriptsize,inner sep=0.8pt] {$2$} (Sh0);
\draw[{Hooks[right]}->] (one) -- node[midway,fill=white,font=\scriptsize,inner sep=0.8pt] {$2$} (P0);
\draw[{Hooks[right]}->] (P0) -- node[pos=0.7,fill=white,font=\scriptsize,inner sep=0.8pt] {$2$} (Sh0);
\draw[-] (one) -- node[pos=0.7,fill=white,font=\scriptsize,inner sep=0.8pt] {$2$} (mt);
\draw[{Hooks[right]}->] (mt) -- node[midway,fill=white,font=\scriptsize,inner sep=0.8pt] {$2$} (D2);
\draw[{Hooks[right]}->] (Sh0) -- node[midway,fill=white,font=\scriptsize,inner sep=0.8pt] {$2$} (Sh);
\draw[{Hooks[right]}->] (one) -- node[midway,fill=white,font=\scriptsize,inner sep=0.8pt] {$2$} (mu);
\draw[{Hooks[right]}->] (P0) -- node[midway,fill=white,font=\scriptsize,inner sep=0.8pt] {$2$} (Sp);
\draw[{Hooks[right]}->] (D2) -- node[midway,fill=white,font=\scriptsize,inner sep=0.8pt] {$2$} (Sh);
\draw[{Hooks[right]}->] (Sh) -- node[midway,fill=white,font=\scriptsize,inner sep=0.8pt] {$6$} (S);
\draw[{Hooks[right]}->] (Spp) -- node[midway,fill=white,font=\scriptsize,inner sep=0.8pt] {$2$} (S);
\draw[-] (mu) -- node[midway,fill=white,font=\scriptsize,inner sep=0.8pt] {$2$} (D2);
\draw[<->,preaction={draw,white,line width=3pt,-}] (Sp) -- node[pos=0.35,fill=white,font=\scriptsize,inner sep=0.8pt] {$2$} (Sh);
\draw[{Hooks[right]}->,preaction={draw,white,line width=3pt,-}] (mu) -- node[pos=0.7,fill=white,font=\scriptsize,inner sep=0.8pt] {$2$} (Sp);
\draw[{Hooks[right]}->,preaction={draw,white,line width=3pt,-}] (Sp) -- node[pos=0.75,fill=white,font=\scriptsize,inner sep=0.8pt] {$6$} (Spp);
\end{tikzpicture}
\end{center}
\legende{Figure de l'édition, redessinée d'après le dessin de la page 133 (dessin 144/07-6). Le cube (II) des groupes $\Pi_0$, $\Pi$, $\Pi'$ et de leurs produits par $\mu$, tous d'indice $2$ les uns dans les autres, puis d'indice $6$ dans $\mathfrak{S}$ et $\underline{\mathfrak{S}}^+$ ; à l'avant les groupes « à $\mu$ près », au fond $\mu_\tau$, $\{1\}$, $\Pi'$, $\Pi_0$. La lecture $\mathfrak{S}^! = \Pi^!_0$ reste incertaine, et les arêtes $\underline{\mu}$ -- $D_2$, $\{1\}$ -- $\mu_\tau$ n'ont pas de pointe lisible.}

\cote{144}{135}
\begin{center}
\begin{tikzpicture}[x=1cm,y=1cm,line cap=round,line join=round,>={Stealth[length=1.6mm,width=1.2mm]},font=\small,every node/.style={inner sep=2pt},]
\node (Sh) at (0,0) {$\Pi'=\mathfrak{S}^!$};
\node (S) at (4.2,0) {$\mathfrak{S}$};
\node (Sph) at (0,-2.8) {$\mathfrak{S}^{+!}$};
\node (Sp) at (4.2,-2.8) {$\mathfrak{S}^+$};
\node (Mh) at (2.4,1.3) {$\mathcal{M}^!$};
\node (M) at (6.6,1.3) {$\mathcal{M}$};
\node (Mph) at (2.4,-1.5) {$\mathcal{M}^{+!}$};
\node (Mp) at (6.6,-1.5) {$\mathcal{M}^+$};
\draw[->] (Mh) -- (M);
\draw[->] (Mph) -- (Mh);
\draw[->] (Mp) -- (M);
\draw[->] (Mph) -- (Mp);
\draw[->] (Sh) -- (Mh);
\draw[->] (S) -- (M);
\draw[->] (Sph) -- (Mph);
\draw[->] (Sp) -- (Mp);
\draw[{Hooks[right]}->] (Sph) -- (Sp);
\draw[<->] (Sph) -- (Sh);
\draw[->,preaction={draw,white,line width=3pt,-}] (Sh) -- (S);
\draw[<->,preaction={draw,white,line width=3pt,-}] (Sp) -- (S);
\end{tikzpicture}
\end{center}
\legende{Figure de l'édition, redessinée d'après le dessin de la page 135 (dessin 144/07-8). Une variante du cube, barrée sur la page par de longues boucles, sur les huit groupes $\mathfrak{S}$, $\mathcal{M}$ avec ou sans $!$ et $+$. Le sommet central, surchargé, est lu $\mathfrak{S}$ ; l'arête verticale entre $\mathfrak{S}$ et $\mathfrak{S}^+$, interrompue sur la page au croisement, est tracée d'un seul trait ; la flèche $\mathfrak{S}^{+!} \to \mathfrak{S}^+$, double dans la transcription, est tracée simple.}

\cote{144}{137}
\begin{center}
\begin{tikzpicture}[x=0.72cm,y=0.8cm,line cap=round,line join=round,>={Stealth[length=1.6mm,width=1.2mm]},font=\small,every node/.style={inner sep=2pt},]
\node (A) at (0,0) {$\mathcal{M}(0)$};
\node (B) at (4.4,0) {$\mathcal{M}(0)\times\mu$};
\node (F) at (0,-3) {$\mu_\tau$};
\node (G) at (4.4,-3) {$D_\tau$};
\node (K) at (0,-6) {$\{1\}$};
\node (mu) at (4.4,-6) {$\underline{\mu}$};
\node (C) at (1.7,-1.05) {$\mathcal{M}^!_0$};
\node (D) at (6.1,-1.05) {$\underline{\mathcal{M}}^!\simeq\mathcal{M}^!_0\times\mu$};
\node (E) at (10.5,-1.05) {$\mathcal{M}$};
\node (L) at (1.7,-4.05) {$\Pi'_0=\mathfrak{S}^!_0$};
\node (Nn) at (6.1,-4.05) {$\Pi'=\mathfrak{S}^!$};
\node (O) at (10.5,-4.05) {$\mathfrak{S}$};
\node (S) at (1.7,-7.05) {$\underline{\Pi}_0$};
\node (T) at (6.1,-7.05) {$\underline{\Pi}=\mathfrak{S}^{+!}$};
\node (U) at (10.5,-7.05) {$\underline{\mathfrak{S}}^+$};
\node (H) at (3.4,-2.1) {$\mathcal{M}^{\Pi_0}_{[0]}$};
\node[label={[font=\scriptsize,inner sep=0pt,xshift=0.25cm]below right:$(\text{ou }\mathcal{M}^![0])$}] (I) at (7.8,-2.1) {$\mathcal{M}^{\Pi}[0]$};
\node[label={[font=\scriptsize,inner sep=1pt]right:$\text{ou }\mathcal{M}[0]$}] (J) at (12.2,-2.1) {$\mathcal{M}^{\mathfrak{S}}[0]$};
\node (Pp) at (3.4,-5.1) {$D^{\Pi_0}_0$};
\node (Q) at (7.8,-5.1) {$D^{\Pi}_0$};
\node (R) at (12.2,-5.1) {$D^{\mathfrak{S}}_0$};
\node (V) at (3.4,-8.1) {$L^{\Pi_0}_0$};
\node (W) at (7.8,-8.1) {$L^{\Pi}_0$};
\node (X) at (12.2,-8.1) {$L^{\mathfrak{S}}_0$};
\draw[{Hooks[right]}->] (H) -- node[midway,fill=white,font=\scriptsize,inner sep=0.8pt] {$2$} (I);
\draw[{Hooks[right]}->] (I) -- node[midway,fill=white,font=\scriptsize,inner sep=0.8pt] {$2$} (J);
\draw[<->] (I) -- (Q);
\draw[<->] (J) -- (R);
\draw[{Hooks[right]}->] (Pp) -- node[midway,fill=white,font=\scriptsize,inner sep=0.8pt] {$2$} (Q);
\draw[{Hooks[right]}->] (Q) -- node[midway,fill=white,font=\scriptsize,inner sep=0.8pt] {$2$} (R);
\draw[{Hooks[right]}->] (V) -- node[pos=0.72,fill=white,font=\scriptsize,inner sep=0.8pt] {$2$} (Pp);
\draw[<->] (W) -- node[pos=0.72,fill=white,font=\scriptsize,inner sep=0.8pt] {$2$} (Q);
\draw[{Hooks[right]}->] (X) -- node[pos=0.3,fill=white,font=\scriptsize,inner sep=0.8pt] {$2$} (R);
\draw[{Hooks[right]}->] (V) -- node[midway,fill=white,font=\scriptsize,inner sep=0.8pt] {$2$} (W);
\draw[{Hooks[right]}->] (W) -- node[midway,fill=white,font=\scriptsize,inner sep=0.8pt] {$2$} (X);
\draw[{Hooks[right]}->] (H) -- (C);
\draw[{Hooks[right]}->] (D) -- (I);
\draw[{Hooks[right]}->] (J) -- (E);
\draw[{Hooks[right]}->] (L) -- (Pp);
\draw[{Hooks[right]}->] (Q) -- (Nn);
\draw[{Hooks[right]}->] (R) -- (O);
\draw[{Hooks[right]}->] (V) -- (S);
\draw[{Hooks[right]}->] (W) -- (T);
\draw[{Hooks[right]}->] (X) -- (U);
\draw[{Hooks[right]}->] (C) -- node[midway,fill=white,font=\scriptsize,inner sep=0.8pt] {$2$} (D);
\draw[{Hooks[right]}->] (D) -- node[midway,fill=white,font=\scriptsize,inner sep=0.8pt] {$6$} (E);
\draw[<->] (C) -- (L);
\draw[{Hooks[right]}->] (L) -- node[midway,fill=white,font=\scriptsize,inner sep=0.8pt] {$2$} (Nn);
\draw[<->] (L) -- node[pos=0.3,fill=white,font=\scriptsize,inner sep=0.8pt] {$2$} (S);
\draw[densely dashed,->] (A) to[out=200,in=180,looseness=1.6] (H);
\draw[densely dashed,->] (F) to[out=200,in=180,looseness=1.6] (Pp);
\draw[densely dashed,->] (K) to[out=200,in=180,looseness=1.6] (V);
\draw[{Hooks[right]}->] (A) -- (C);
\draw[{Hooks[right]}->] (B) -- (D);
\draw[{Hooks[right]}->] (F) -- (L);
\draw[{Hooks[right]}->] (G) -- (Nn);
\draw[{Hooks[right]}->] (K) -- node[midway,fill=white,font=\scriptsize,inner sep=0.8pt] {$2$} (L);
\draw[{Hooks[right]}->] (L) to[out=-75,in=180] (mu);
\draw[{Hooks[right]}->] (K) -- (S);
\draw[{Hooks[right]}->] (mu) -- (T);
\draw[{Hooks[right]}->] (A) -- node[midway,fill=white,font=\scriptsize,inner sep=0.8pt] {$2$} (B);
\draw[<->] (A) -- (F);
\draw[<->] (F) -- node[pos=0.3,fill=white,font=\scriptsize,inner sep=0.8pt] {$2$} (K);
\draw[{Hooks[right]}->,preaction={draw,white,line width=3pt,-}] (Nn) -- (D);
\draw[{Hooks[right]}->,preaction={draw,white,line width=3pt,-}] (O) -- (E);
\draw[{Hooks[right]}->,preaction={draw,white,line width=3pt,-}] (Nn) -- node[midway,fill=white,font=\scriptsize,inner sep=0.8pt] {$6$} (O);
\draw[{Hooks[right]}->,preaction={draw,white,line width=3pt,-}] (T) -- node[pos=0.3,fill=white,font=\scriptsize,inner sep=0.8pt] {$2$} (Nn);
\draw[{Hooks[right]}->,preaction={draw,white,line width=3pt,-}] (U) -- node[pos=0.3,fill=white,font=\scriptsize,inner sep=0.8pt] {$2$} (O);
\draw[{Hooks[right]}->,preaction={draw,white,line width=3pt,-}] (S) -- node[midway,fill=white,font=\scriptsize,inner sep=0.8pt] {$2$} (T);
\draw[{Hooks[right]}->,preaction={draw,white,line width=3pt,-}] (T) -- node[midway,fill=white,font=\scriptsize,inner sep=0.8pt] {$6$} (U);
\draw[{Hooks[right]}->,preaction={draw,white,line width=3pt,-}] (G) -- (B);
\draw[{Hooks[right]}->,preaction={draw,white,line width=3pt,-}] (F) -- node[midway,fill=white,font=\scriptsize,inner sep=0.8pt] {$2$} (G);
\end{tikzpicture}
\end{center}
\legende{Figure de l'édition, redessinée d'après le dessin de la page 137 (dessin 144/07-9). Le diagramme (III), sur la place de $\mathcal{M}(0)$ : trois couches de cubes accolés, à l'avant $\mathcal{M}(0)$, $\mu_\tau$, $\{1\}$ et les groupes au-dessus d'eux, au milieu les chaînes $\mathcal{M}^!_0 \to \underline{\mathcal{M}}^! \to \mathcal{M}$, $\Pi'_0 \to \Pi' \to \mathfrak{S}$, $\underline{\Pi}_0 \to \underline{\Pi} \to \underline{\mathfrak{S}}^+$, au fond les chaînes en $[0]$, $D_0$ et $L_0$ ; les arcs pointillés relient $\mathcal{M}(0)$, $\mu_\tau$, $\{1\}$ au fond. Les arêtes sont disposées en cube d'après le dessin ; au centre, où les traits se croisent sous les étiquettes, les sommets que joignent quelques courtes flèches (dont celle qui aboutit à $\underline{\mu}$) sont une lecture.}

\cote{144}{139}
\begin{center}
\begin{tikzpicture}[x=1.25mm,y=1.25mm,line cap=round,line join=round,>={Stealth[length=1.6mm,width=1.2mm]},font=\small,every node/.style={inner sep=2pt},]
\node[circle,draw,inner sep=1.3pt] (X1) at (22,-12.5) {};
\node (a) at (41,-12.5) {$N^?_\rho$};
\node[circle,draw,inner sep=1.3pt] (X2) at (33,-18) {};
\node (b) at (50,-17.5) {$N^?_\rho$};
\node (c) at (67.5,-17.5) {$\mathcal{M}^?$};
\node (d) at (88.5,-17.5) {$\mathcal{M}$};
\node (e) at (9.7,-24.5) {$\mathcal{M}^?_0$};
\node (f) at (29.4,-24.5) {$\mathcal{M}^!$};
\node (g) at (17,-30) {$Z^!_{0\rho}=\mathcal{M}(\bar{\jmath})$};
\node (h) at (41,-29.5) {$Z^!_\rho$};
\node (i) at (59,-29.5) {$\mathcal{M}'^!$};
\node (j) at (80.5,-29.5) {$\mathcal{M}'$};
\node (k) at (34.8,-36) {$\mu_{\tau'}$};
\node (l) at (50,-36) {$D_{\tau'}$};
\node (m) at (67.5,-36) {$\Pi_{\tau'}$};
\node (n) at (88.5,-36) {$\mathfrak{S}$};
\node (o) at (21,-38) {$\mathcal{M}$};
\node (p) at (10.2,-44) {$\mathcal{M}'^!_0$};
\node (q) at (30.3,-44) {$\mathcal{M}'^!$};
\node (r) at (21.8,-49.5) {$1$};
\node (s) at (40,-50.6) {$\underline{\mu}$};
\node (t) at (58.5,-50.6) {$\underline{\Pi}$};
\node (u) at (80,-50.6) {$\underline{\mathfrak{S}}^+$};
\node (v) at (21,-59.6) {$\mathcal{M}'$};
\node (w) at (10.6,-64) {$\underline{\Pi}_0$};
\node (x) at (30.8,-64) {$\underline{\Pi}$};
\node (y) at (21.8,-78.4) {$\mathfrak{S}^+$};
\draw[{Hooks[right]}->] (b) -- (c);
\draw[{Hooks[right]}->] (c) -- node[midway,fill=white,font=\scriptsize,inner sep=0.8pt] {$6$} (d);
\draw[{Hooks[right]}->] (l) -- (b);
\draw[{Hooks[right]}->] (m) -- (c);
\draw[{Hooks[right]}->] (n) -- (d);
\draw[{Hooks[right]}->] (l) -- (m);
\draw[{Hooks[right]}->] (m) -- node[midway,fill=white,font=\scriptsize,inner sep=0.8pt] {$6$} (n);
\draw[densely dashed,->] (b) to[bend right=40] (a);
\draw[{Hooks[right]}->] (X1) -- node[midway,fill=white,font=\scriptsize,inner sep=0.8pt] {$2$} (a);
\draw[{Hooks[right]}->] (a) -- (f);
\draw[{Hooks[right]}->] (X2) -- (b);
\draw[{Hooks[right]}->] (k) -- node[midway,fill=white,font=\scriptsize,inner sep=0.8pt] {$2$} (l);
\draw[{Hooks[right]}->] (X1) -- (e);
\draw[{Hooks[right]}->] (g.north -| X1) -- node[midway,fill=white,font=\scriptsize,inner sep=0.8pt] {$2$} (X1);
\draw[{Hooks[right]}->] (f) -- node[midway,fill=white,font=\scriptsize,inner sep=0.8pt] {$2$} (X2);
\draw[{Hooks[right]}->] (e) -- node[midway,fill=white,font=\scriptsize,inner sep=0.8pt] {$2$} (f);
\draw[{Hooks[right]}->] (h) -- node[midway,fill=white,font=\scriptsize,inner sep=0.8pt] {$2$} (b);
\draw[{Hooks[right]}->] (i) -- node[midway,fill=white,font=\scriptsize,inner sep=0.8pt] {$2$} (c);
\draw[{Hooks[right]}->] (j) -- node[midway,fill=white,font=\scriptsize,inner sep=0.8pt] {$2$} (d);
\draw[{Hooks[right]}->] (s) -- node[midway,fill=white,font=\scriptsize,inner sep=0.8pt] {$2$} (l);
\draw[{Hooks[right]}->] (t) -- node[midway,fill=white,font=\scriptsize,inner sep=0.8pt] {$2$} (m);
\draw[{Hooks[right]}->] (u) -- node[midway,fill=white,font=\scriptsize,inner sep=0.8pt] {$2$} (n);
\draw[{Hooks[right]}->] (q) -- node[midway,fill=white,font=\scriptsize,inner sep=0.8pt] {$2$} (f);
\draw[{Hooks[right]}->] (o) -- (g);
\draw[{Hooks[right]}->] (f) -- node[pos=0.75,fill=white,font=\scriptsize,inner sep=0.8pt] {$6$} (o);
\draw[{Hooks[right]}->] (p) -- node[midway,fill=white,font=\scriptsize,inner sep=0.8pt] {$2$} (e);
\draw[{Hooks[right]}->] (p) -- node[midway,fill=white,font=\scriptsize,inner sep=0.8pt] {$2$} (q);
\draw[{Hooks[right]}->] (g) -- (p);
\draw[{Hooks[right]}->] (v) -- node[pos=0.25,fill=white,font=\scriptsize,inner sep=0.8pt] {$2$} (o);
\draw[{Hooks[right]}->] (x) -- (q);
\draw[{Hooks[right]}->] (q) -- node[midway,fill=white,font=\scriptsize,inner sep=0.8pt] {$6$} (v);
\draw[{Hooks[right]}->] (w) -- (p);
\draw[{Hooks[right]}->] (r) -- (w);
\draw[{Hooks[right]}->] (w) -- node[midway,fill=white,font=\scriptsize,inner sep=0.8pt] {$2$} (x);
\draw[{Hooks[right]}->] (y) -- (v);
\draw[{Hooks[right]}->] (x) -- node[midway,fill=white,font=\scriptsize,inner sep=0.8pt] {$6$} (y);
\draw[{Hooks[right]}->] (k) -- (h);
\draw[{Hooks[right]}->] (s) -- (x);
\draw[densely dashed,<->] (t) to[out=-120,in=0] (x);
\draw[densely dashed,<->] (u) to[out=-110,in=0] (y);
\draw[{Hooks[right]}->] (s) -- (t);
\draw[{Hooks[right]}->] (t) -- node[midway,fill=white,font=\scriptsize,inner sep=0.8pt] {$6$} (u);
\draw[{Hooks[right]}->,preaction={draw,white,line width=3pt,-}] (g) -- node[midway,fill=white,font=\scriptsize,inner sep=0.8pt] {$2$} (h);
\draw[{Hooks[right]}->,preaction={draw,white,line width=3pt,-}] (r) -- node[midway,fill=white,font=\scriptsize,inner sep=0.8pt] {$2$} (s);
\draw[{Hooks[right]}->,preaction={draw,white,line width=3pt,-}] (h) -- (i);
\draw[{Hooks[right]}->,preaction={draw,white,line width=3pt,-}] (i) -- node[midway,fill=white,font=\scriptsize,inner sep=0.8pt] {$6$} (j);
\draw[{Hooks[right]}->,preaction={draw,white,line width=3pt,-}] (s) -- (h);
\draw[{Hooks[right]}->,preaction={draw,white,line width=3pt,-}] (t) -- (i);
\draw[{Hooks[right]}->,preaction={draw,white,line width=3pt,-}] (u) -- (j);
\draw[densely dashed,->] (29.4,-11.2) -- (25.4,-11.2);
\draw[{Hooks[right]}->] (25.9,-42.6) -- (29.6,-38.6);
\end{tikzpicture}
\end{center}
\legende{Figure de l'édition, redessinée d'après le dessin de la page 139 (dessin 144/07-11), mise au net du diagramme (IV) de la page 137 (dessin 144/07-10), dont elle suit la disposition en trois blocs de cubes : centralisateurs $Z_\rho$, normalisateurs $N_\rho$, groupes $\mathcal{M}$, $\mathcal{M}'$, au-dessus de $\mu_{\tau'}$, $D_{\tau'}$, $\Pi_{\tau'}$, $\mathfrak{S}$ et de $\underline{\mu}$, $\underline{\Pi}$, $\underline{\mathfrak{S}}^+$. Les deux sommets barrés sur la page ($N^?_{0\rho}$ et l'un des $N^?_\rho$) sont des ronds vides ; deux courtes flèches sans sommet (en pointillé en haut, et dans le bloc de gauche) sont laissées telles qu'elles sont tracées ; dans le bloc de gauche, les sommets que joignent quelques traits courts sont une lecture. L'esquisse au crayon du bas de la page n'est pas reprise.}

\cote{156-4}{51}
\begin{center}
\begin{tikzpicture}[x=1.1cm,y=1cm,line cap=round,line join=round,>={Stealth[length=1.6mm,width=1.2mm]},font=\small,every node/.style={inner sep=2pt},]
\node (Phi) at (-1.9,0.45) {$\Phi\subset$};
\node (Fb) at (0,0) {$\overline{F}'$};
\node (Ft) at (1.9,1.3) {$\widetilde{F}$};
\node (Fp) at (1.1,-0.95) {$F'$};
\node (Fv) at (3.7,0.45) {$\vec{F}$};
\node (F) at (5.6,0) {$F$};
\node (Gb) at (0,-3.2) {$\overline{G}'$};
\node (Gt) at (1.9,-1.9) {$\widetilde{G}$};
\node (Gp) at (1.1,-4.15) {$G'$};
\node (Gv) at (3.7,-2.75) {$\overline{G}$};
\node (G) at (5.6,-3.2) {$G$};
\draw[-] (Ft) -- (Gt);
\draw[{Hooks[right]}->] (Fb) to[bend left=15] (Ft);
\draw[->] (Ft) -- node[midway,above,sloped,font=\scriptsize,inner sep=1.5pt] {$\sim$} (Fv);
\draw[->] (Gt) -- node[midway,above,sloped,font=\scriptsize,inner sep=1.5pt] {$\sim$} (Gv);
\draw[->] (Phi) -- (Fb);
\draw[->] (Fv) -- node[midway,above,sloped,font=\scriptsize,inner sep=1.5pt] {$\sim$} (F);
\draw[->] (Gv) -- node[midway,above,sloped,font=\scriptsize,inner sep=1.5pt] {$\sim$} (G);
\draw[-] (Fb) -- (Gb);
\draw[-] (Fv) -- (Gv);
\draw[-] (F) -- (G);
\draw[->] (Fb) -- node[midway,above,sloped,font=\scriptsize,inner sep=1.5pt] {$\sim$} (Fp);
\draw[->] (Gb) -- node[midway,above,sloped,font=\scriptsize,inner sep=1.5pt] {$\sim$} (Gp);
\draw[-] (Fp) -- (Gp);
\draw[{Hooks[right]}->,preaction={draw,white,line width=3pt,-}] (Fp) to[bend left=10] (Fv);
\draw[{Hooks[right]}-] (Gb) -- ($(Gb)!0.5!(Gt)$);
\draw[{Hooks[right]}-] (Gp) -- ($(Gp)!0.5!(Gv)$);
\end{tikzpicture}
\end{center}
\legende{Figure de l'édition, redessinée d'après le dessin de la page 51 (dessin 156-4/03-5). Le cube de la démonstration : en haut $\Phi \subset \overline{F}'$, $\widetilde{F}$, $\vec{F}$, $F$ et $F'$, en bas $\overline{G}'$, $\widetilde{G}$, $\overline{G}$, $G$ et $G'$ ; $\sim$ marque les subdivisions, les flèches à crochet les inclusions, et les traits verticaux joignent chaque figure du haut à celle du bas qui lui répond. Le trait de $\widetilde{F}$ à $\widetilde{G}$, interrompu sur la page à la hauteur de $F'$, est lu comme une seule arête passant derrière $F'$ ; les deux crochets d'inclusion du bas restent amorcés comme sur la page. Le grand arc qui entoure le dessin n'est pas repris.}

\cote{156-6}{74}
\begin{center}
\begin{tikzpicture}[x=0.95cm,y=0.95cm,line cap=round,line join=round,>={Stealth[length=1.6mm,width=1.2mm]},font=\small,every node/.style={inner sep=2pt},]
\node (HK) at (0,0) {$H'_K$};
\node (HF) at (-2.7,0.65) {$H'_F$};
\node (HG) at (1.9,1.15) {$H'_G$};
\node (Hp) at (-0.8,1.8) {$H'$};
\node (K) at (0,3) {$K$};
\node (F) at (-2.7,3.65) {$F$};
\node (G) at (1.9,4.15) {$G$};
\node (H) at (-0.8,4.8) {$H$};
\node (FK) at (-1.9,-1.15) {$F'_K$};
\node (GK) at (2.7,-0.65) {$G'_K$};
\node (Kp) at (0.8,-1.8) {$K'$};
\node (Fp) at (-4.6,-0.5) {$F'$};
\node (Gp) at (4.6,0.5) {$G'$};
\draw[->] (Hp) -- (H);
\draw[{Hooks[right]}->] (G) -- (H);
\draw[{Hooks[right]}->] (HG) -- (Hp);
\draw[{Hooks[right]}->] (F) -- (H);
\draw[{Hooks[right]}->] (HF) -- (Hp);
\draw[->] (HG) -- (G);
\draw[->] (HF) -- (F);
\draw[{Hooks[right]}->] (Gp) -- (HG);
\draw[{Hooks[right]}->] (K) -- (G);
\draw[{Hooks[right]}->] (HK) -- (HG);
\draw[{Hooks[right]}->] (HK) -- (HF);
\draw[{Hooks[right]}->] (Fp) -- (HF);
\draw[{Hooks[right]}->] (GK) -- (Gp);
\draw[{Hooks[right]}->] (GK) -- (HK);
\draw[{Hooks[right]}->] (FK) -- (HK);
\draw[{Hooks[right]}->] (FK) -- (Fp);
\draw[{Hooks[right]}->] (Kp) -- (GK);
\draw[{Hooks[right]}->] (Kp) -- (FK);
\draw[{Hooks[right]}->,preaction={draw,white,line width=3pt,-}] (K) -- (F);
\draw[->,preaction={draw,white,line width=3pt,-}] (HK) -- (K);
\end{tikzpicture}
\end{center}
\legende{Figure de l'édition, redessinée d'après le dessin de la page 74 (dessin 156-6/04-4). Le cube tronqué entier : en haut le carré $K$, $F$, $G$, $H$ ; sous lui le carré $H'_K$, $H'_F$, $H'_G$, $H'$, relié au premier par les flèches verticales, qui sont des $\ll$ ; en bas $F'$, $G'$, $F'_K$, $G'_K$ et $K'$. Les flèches à crochet sont des $\leq$.}

\cote{156-9}{112}
\begin{center}
\begin{tikzpicture}[x=0.9cm,y=0.9cm,line cap=round,line join=round,>={Stealth[length=1.6mm,width=1.2mm]},font=\small,every node/.style={inner sep=2pt},]
\node (o) at (0,0) {$0$};
\node (a) at (-2.1,2.1) {$a$};
\node (c) at (0.6,1.5) {$c$};
\node (bp) at (2.3,2.1) {$b'$};
\node (B) at (-1.5,3.6) {$B$};
\node (C) at (0.2,4.2) {$C$};
\node (A) at (2.9,3.6) {$A$};
\node (one) at (0.8,5.7) {$1$};
\node (Ap) at (1.85,4.65) {$A'$};
\draw[->] (c) -- (B);
\draw[->] (c) -- (A);
\draw[->] (B) -- (one);
\draw[->] (A) -- (Ap);
\draw[->] (Ap) -- (one);
\draw[->] (o) -- (c);
\draw[->] (a) -- (B);
\draw[->] (bp) -- (A);
\draw[->] (C) -- (one);
\draw[->] (o) -- node[midway,below left,font=\scriptsize,inner sep=1.5pt] {$a'$} (a);
\draw[->] (o) -- (bp);
\draw[->,preaction={draw,white,line width=3pt,-}] (a) -- (C);
\draw[->,preaction={draw,white,line width=3pt,-}] (bp) -- (C);
\node at (0.75,2.60) {$\ast$};
\end{tikzpicture}
\end{center}
\legende{Figure de l'édition, redessinée d'après le dessin de la page 112 (dessin 156-9/06-1). Le treillis de la page, dessiné comme un cube vu de biais : de $0$ partent des flèches vers $a$, $c$ et $b'$ ($a'$ est écrit le long de la flèche vers $a$), de là vers $B$, $C$, $A$, puis vers $1$, en passant par $A'$ entre $A$ et $1$ ; l'étoile du centre est celle de la page. Les lettres du dessin ne sont pas toutes celles du texte, et l'hexagone orienté dessiné au-dessous n'est pas repris.}


\section{Ce qu'on peut en tirer}

Deux groupes de dessins portent des énoncés que l'on pourrait vérifier en Lean. Les graphes stables de la cote 147, rangés par type $(g, \nu)$ avec la dimension $3g - 3 + \nu$, se comptent ; c'est de la combinatoire finie. Les fibres du revêtement $x \mapsto x^n$ au-dessus de $0$, $1$, $\infty$ (cote 138) se calculent sur des polynômes explicites. Les figures de la cote 63 (pages 97 et 99) et de la cote 84 (page 125) illustrent des énoncés déjà démontrés dans l'article Lean ; elles pourraient y entrer sous la forme d'une figure de l'édition tirée des énoncés eux-mêmes, plutôt que d'une copie de la page. Le reste relève de l'histoire des mathématiques : la façon dont Grothendieck se représentait $\operatorname{Spec}\mathbf Z$, une famille, une courbe stable.

\section*{Déclaration sur l'usage de l'intelligence artificielle}
Les figures ont été redessinées par Claude (Opus 5.5, Anthropic) à partir des dessins fidèles de l'édition, et ce brouillon a été assemblé par le même modèle à partir des lectures modernisées. Aucune figure n'a été comparée par un humain au fac-similé.
\end{document}
